\documentclass[12pt, a4paper]{article}

\usepackage[utf8]{inputenc}
\usepackage[T1]{fontenc}
\usepackage{lmodern}
\usepackage{fix-cm}

\usepackage[spanish]{babel}
\addto\captionsspanish{\renewcommand{\tablename}{Tabla}}

\usepackage{amsmath, amssymb}
\usepackage{graphicx}
\usepackage{xcolor}

\usepackage{pgfplots}
\pgfplotsset{compat=1.18}

\usepackage{tikz}
\usetikzlibrary{shapes.geometric, arrows.meta, positioning}
\usepackage{geometry}
\usepackage{booktabs}
\usepackage{array}
\usepackage{fancyhdr}
\usepackage{float}
\usepackage{caption}
\usepackage{listings}
\usepackage{pdflscape}
\usepackage{multirow}
\usepackage{enumitem}
\usepackage{hanging}
\usepackage{xurl}
\usepackage{hyperref}
\hypersetup{hidelinks}

\geometry{margin=2cm, top=3cm, bottom=3cm}
\setlength{\headheight}{13.6pt}

% Configuración global de tablas
\renewcommand{\arraystretch}{1.2}
\captionsetup[table]{skip=8pt}

% Tipos de columna con salto de línea automático
\newcolumntype{L}[1]{>{\raggedright\arraybackslash}p{#1}}   % Izquierda con ancho
\newcolumntype{R}[1]{>{\raggedleft\arraybackslash}p{#1}}    % Derecha con ancho
\newcolumntype{C}[1]{>{\centering\arraybackslash}p{#1}}     % Centrado con ancho
\newcolumntype{J}[1]{>{\arraybackslash}p{#1}}               % Justificado con ancho

% Para tabularx (ancho flexible)
\newcolumntype{X}{>{\raggedright\arraybackslash}X}
\newcolumntype{Y}{>{\centering\arraybackslash}X}
\newcolumntype{Z}{>{\raggedleft\arraybackslash}X}

% Entorno de tabla estándar
\newenvironment{mitabla}[2][H]{%
  \begin{table}[#1]
  \centering
  \caption{#2}
  \small
}{%
  \end{table}
}

\pagestyle{fancy}
\fancyhf{}
\rhead{\small Análisis de Mercado y Herramienta de Recolección de Datos}
\lhead{\small Informe de Actividad}
\cfoot{\thepage}
\renewcommand{\headrulewidth}{0.4pt}

\lstset{
  basicstyle=\ttfamily\small,
  keywordstyle=\bfseries,
  commentstyle=\itshape,
  stringstyle=\ttfamily,
  numbers=left,
  numberstyle=\tiny,
  frame=single,
  breaklines=true,
  tabsize=4,
  literate=
    {á}{{\'a}}1 {é}{{\'e}}1 {í}{{\'\i}}1 {ó}{{\'o}}1 {ú}{{\'u}}1
    {Á}{{\'A}}1 {É}{{\'E}}1 {Í}{{\'I}}1 {Ó}{{\'O}}1 {Ú}{{\'U}}1
    {ñ}{{\~n}}1 {Ñ}{{\~N}}1
}

% Marca visible (en negrita) para resultados calculados con el dataset de prueba.
% Cuando se procesen las respuestas reales: reemplazar los valores y
% redefinir estos comandos como vacíos.
\newcommand{\datosprueba}{\textbf{[Datos de prueba]}}
\newcommand{\avisoprueba}{\begin{center}\fbox{\parbox{0.92\textwidth}{\small\textbf{Resultados provisionales.} Los valores de esta sección se calcularon con el dataset de prueba (respuestas sintéticas) para validar la herramienta. Se reemplazan por los resultados de las encuestas reales antes de la entrega.}}\end{center}}

\begin{document}

\setlength{\parindent}{0pt}
\pagenumbering{roman}

\begin{titlepage}
\centering
{\large\bfseries UNIVERSIDAD MAYOR DE SAN SIMÓN\par}
\vspace{4pt}
{\bfseries FACULTAD DE CIENCIAS Y TECNOLOGÍA\par}
\vspace{4pt}
{\bfseries CARRERA DE INGENIERÍA DE SISTEMAS\par}

\vspace*{\fill}
{\huge\bfseries ANÁLISIS DE MERCADO Y HERRAMIENTA\\DE RECOLECCIÓN DE DATOS\par}
\vspace{0.8cm}
{\Large Enfoque ARCUSUR\par}
\vspace{0.8cm}
{\normalsize Informe de Actividad 2\par}
\vspace*{\fill}

\begin{flushleft}
\begin{tabular}{@{}ll}
\textbf{Presentado por:} & Miguel Mita Choque \\[4pt]
\textbf{Docente:}        & Henrry Frank Villarroel Tapia \\[4pt]
\textbf{Materia:}        & Taller de Simulación de Sistemas \\
\end{tabular}
\end{flushleft}

\vspace{1.5cm}
{\bfseries Cochabamba, Bolivia\par}
\end{titlepage}

\section*{Resumen}

Este informe presenta el análisis de mercado sobre el perfil, la
formación continua y la empleabilidad de los titulados de la Carrera de
Ingeniería de Sistemas de la Universidad Mayor de San Simón, como insumo
para el proceso de Autoevaluación orientado a la acreditación ARCU-SUR.
Las variables analizadas se eligieron a partir de los indicadores que el
documento de Criterios de Calidad ARCU-SUR para Ingeniería pide sobre
graduados, perfil de egreso, posgrado y vinculación con el sector
productivo, y se obtienen de dos encuestas: una a titulados y otra a
empleadores. Los datos se procesan en una herramienta de software
propia que calcula frecuencias, medias con intervalo de confianza del
95\%, pruebas de independencia chi-cuadrada sobre tablas de
contingencia, y simula escenarios de inserción laboral mediante el
método de la Transformada Inversa discreta. \textit{[Pendiente: completar
con la síntesis de resultados una vez procesadas las respuestas reales.]}

\textbf{Palabras clave:} análisis de mercado, empleabilidad, ARCU-SUR,
simulación, Ingeniería de Sistemas.

\tableofcontents
\listoftables
\newpage

\pagenumbering{arabic}
\setlength{\parskip}{10pt}

\section{Introducción y Justificación}

\subsection{Antecedentes}

La Carrera de Ingeniería de Sistemas de la Universidad Mayor de San Simón
se encuentra en un proceso de Autoevaluación orientado a la elaboración de
un Plan de Desarrollo Estratégico, en el marco de una futura postulación
de acreditación ante el sistema ARCU-SUR. Este proceso requiere contar con
información empírica actualizada sobre el perfil, las expectativas de
formación continua y la situación de empleabilidad de sus titulados, así
como sobre los requerimientos que el sector productivo plantea a los
profesionales egresados de la carrera.

El propio documento de Criterios de Calidad ARCU-SUR, en su apartado de
documentación y fuentes de información, pide presentar las ``encuestas
aplicadas para la autoevaluación'' y la ``información sobre mecanismos
de seguimiento de graduados'' (ARCU-SUR, 2019). El presente estudio
levanta esa información mediante dos instrumentos propios (encuestas a
titulados y a empleadores) y la procesa con una herramienta de software
desarrollada para este fin.

\subsection{Objetivos}

\subsubsection{Objetivo General}

Analizar el perfil, la formación continua y la empleabilidad de los
titulados de la Carrera de Ingeniería de Sistemas de la UMSS, de acuerdo
con los criterios de calidad ARCU-SUR para graduados y perfil de egreso.

\subsubsection{Objetivos Específicos}

\begin{enumerate}
  \item Calcular los indicadores de condiciones de empleo que pide el
  criterio ARCU-SUR 3.2.3 (tasa de empleo, tiempo medio para obtener el
  primer empleo y concordancia entre titulación y empleo).
  \item Determinar, mediante pruebas de independencia chi-cuadrada, la
  relación entre la situación laboral, la formación complementaria y la
  valoración de la formación recibida.
  \item Describir la valoración que los empleadores hacen del perfil
  profesional y de las competencias de los titulados.
  \item Simular escenarios de inserción laboral a partir de las
  distribuciones observadas en la encuesta.
  \item Desarrollar una herramienta de software que importe, procese y
  visualice los datos recolectados.
\end{enumerate}

\subsection{Justificación}

El sistema ARCU-SUR evalúa a las carreras de ingeniería en cuatro
dimensiones: Contexto Institucional, Proyecto Académico, Comunidad
Universitaria e Infraestructura (ARCU-SUR, 2019). Varios de sus
criterios dependen directamente de información sobre los titulados y
sobre quienes los emplean:

\begin{itemize}
  \item \textbf{3.2.3 Condiciones de empleo}, cuyos indicadores son el
  ``tiempo medio para obtener el primer empleo'', el ``destino laboral y
  profesional de los graduados; tasa de empleo'' y la ``concordancia
  entre las características de la titulación y las del empleo''.
  \item \textbf{3.2.1 Resultados} y \textbf{3.2.2 Vinculación y
  seguimiento a los graduados}, que piden evaluar el resultado del
  proceso formativo y contar con un sistema de seguimiento de
  graduados.
  \item \textbf{2.1.2 Perfil de Egreso}, que exige ``coherencia entre el
  perfil de egreso y la demanda explícita de competencias profesionales
  \ldots{} expresadas por agentes sociales relevantes''.
  \item \textbf{2.1.7 Actualización curricular}, \textbf{2.4.1 Cursos de
  actualización profesional permanente} y \textbf{1.1.5} (programas de
  posgrado).
\end{itemize}

Este estudio produce, con datos propios, la evidencia que esos criterios
piden. Además, la herramienta de software queda como un mecanismo
permanente de seguimiento: cada nueva aplicación de las encuestas se
importa y los indicadores se recalculan sin rehacer el análisis.

\section{Marco Metodológico}

\subsection{Diseño}

El estudio es de tipo cuantitativo y exploratorio. Es cuantitativo
porque se basa en respuestas estructuradas (categóricas y numéricas)
resumidas mediante frecuencias, medias e intervalos de confianza. Es
exploratorio porque busca identificar relaciones entre variables sin
partir de una hipótesis previa. Se complementa con un componente de
simulación: a partir de las distribuciones observadas se generan
cohortes de titulados para estudiar cómo cambiarían los indicadores de
empleo bajo distintos escenarios.

\subsection{Población Objetivo}

\begin{enumerate}
  \item \textbf{Titulados} de la Carrera de Ingeniería de Sistemas de la
  UMSS. La encuesta se distribuye únicamente a titulados de la Carrera,
  por lo que el formulario no incluye una pregunta filtro.
  \item \textbf{Empleadores} de titulados de la Carrera: cinco
  organizaciones (cantidad confirmada con el docente), con al menos una
  del ámbito público y una del privado, como exige el enunciado.
\end{enumerate}

\subsection{Técnicas de Recolección de Datos}

Se aplican dos encuestas estructuradas mediante formularios web: la de
titulados en Microsoft Forms y la de empleadores en Google Forms. Ambos
instrumentos se transcriben completos en los Anexos del enunciado. La
encuesta a titulados tiene lógica condicional: según la situación
laboral declarada, se muestra una rama distinta de preguntas
(trabaja en una organización, tiene un emprendimiento propio o no
trabaja).

Ambos formularios registran automáticamente la fecha de cada respuesta.
La herramienta usa esa fecha para delimitar el período analizado, de
modo que el formulario puede seguir abierto y cada informe se cierra con
una fecha de corte, o se comparan respuestas antiguas y nuevas.

% Generado por sistema/backend ReporteTex.java -- no editar a mano.
% Datos: Base de datos: 148 importado (titulados) / Base de datos: 53 importado (empleadores)
Se analizan 148 respuestas de titulados, recibidas entre el 29 de abril de 2026 y el 18 de octubre de 2026, y 53 respuestas de empleadores, recibidas entre el 11 de septiembre de 2026 y el 17 de octubre de 2026. La fecha de cada respuesta proviene del registro automático del formulario.



Las Historias de Usuario del enunciado orientaron el diseño de la
herramienta. La Tabla~\ref{tab:hu} muestra qué preguntas reales del
formulario responden a cada una; el análisis se hace siempre sobre las
preguntas tal como fueron aplicadas.

\begin{mitabla}{Historias de Usuario y preguntas del formulario que las responden\label{tab:hu}}
\begin{tabular}{cL{4.2cm}L{8.8cm}}
\toprule
HU & Tema & Preguntas del formulario utilizadas \\
\midrule
1 & Antigüedad y empleabilidad & Años de vida profesional, tiempo desempleado, situación laboral, sector, rubro, cargo, remuneración, tiempo hasta el primer empleo \\
2 & Formación continua y áreas de interés & Formación complementaria y su nivel, interés en posgrado, área de interés, institución elegida, opinión sobre la oferta de posgrado de la FCyT \\
3 & Mecanismos de financiamiento & Fuente de financiamiento del programa cursado; financiamiento del emprendimiento \\
4 & Brechas en competencias & Satisfacción con la formación, concordancia formación--mercado, aspectos útiles y a mejorar, asignaturas; escala de perfil profesional de empleadores \\
5 & Reportes y matrices cruzadas & Exportación CSV, tablas de contingencia, mapas de calor y diagrama de radar en la herramienta \\
\bottomrule
\end{tabular}
\end{mitabla}

\subsection{Variables e Indicadores según ARCU-SUR}

Una variable o un cruce se incluye en el análisis solo si responde a un
indicador concreto del documento de Criterios ARCU-SUR y sale de una
pregunta real del formulario. La Tabla~\ref{tab:indicadores} resume los
indicadores calculados y la Tabla~\ref{tab:cruces} los cruces de
variables.

\begin{mitabla}{Indicadores calculados y criterio ARCU-SUR al que responden\label{tab:indicadores}}
\begin{tabular}{cL{4.6cm}L{6.6cm}c}
\toprule
Id & Indicador & Cálculo & Criterio \\
\midrule
U1 & Tasa de empleo & Proporción que trabaja en una organización o en un emprendimiento propio & 3.2.3 \\
U2 & Tiempo medio al primer empleo & Media e IC 95\% usando el punto medio de cada rango (meses) & 3.2.3 \\
U3 & Años de vida profesional y desempleo & Media, desviación estándar e IC 95\% & 3.2.3 \\
U4 & Destino laboral & Frecuencias de rubro, sector, departamento, cargo y remuneración & 3.2.3 \\
U5 & Concordancia titulación--empleo & Frecuencias y media de las escalas de pertinencia y concordancia (1--4) & 3.2.3 \\
U6 & Satisfacción con la formación & Frecuencias y media; aspectos útiles y a mejorar & 3.2.1 \\
U7 & Ranking de asignaturas & Menciones de asignaturas con ventaja y poco útiles & 2.1.7 \\
U8 & Demanda de posgrado & \% con formación complementaria, \% con interés, área e institución & 2.4.1 / 1.1.5 \\
U9 & Oferta de posgrado FCyT & Frecuencias y media de la escala (1--4) & 2.4.1 \\
U10 & Seguimiento de graduados & \% interesado en la red de contactos de la Carrera & 3.2.2 \\
U11 & Perfil según empleadores & Porcentaje de cada opción en los 10 criterios de perfil profesional & 2.1.2 / 2.4.2 \\
\bottomrule
\end{tabular}
\end{mitabla}

\begin{mitabla}{Cruces de variables (pruebas de independencia)\label{tab:cruces}}
\begin{tabular}{cL{5.4cm}L{5.8cm}c}
\toprule
Id & Cruce & Propósito & Criterio \\
\midrule
C1 & Situación laboral $\times$ satisfacción con la formación & Determinar si la valoración de la formación cambia según la situación laboral & 3.2.1 / 3.2.3 \\
C2 & Sector $\times$ pertinencia del puesto & Determinar si la concordancia titulación--empleo depende del sector & 3.2.3 \\
C3 & Formación complementaria $\times$ remuneración & Determinar si la formación complementaria se asocia con la remuneración & 2.4.1 / 3.2.3 \\
C4 & Interés en posgrado $\times$ oferta FCyT afín & Determinar si la afinidad percibida de la oferta FCyT depende del interés en un posgrado & 2.4.1 / 1.1.5 \\
C5 & Área de interés $\times$ institución elegida & Identificar áreas en que la oferta de la UMSS no capta la demanda de posgrado & 2.4.1 / 1.1.5 \\
C6 & Concordancia $\times$ tiempo al primer empleo & Determinar si una formación alineada al mercado se asocia con un ingreso más rápido al empleo & 3.2.3 \\
C7 & Financiamiento $\times$ nivel cursado & Complementar la historia de usuario de financiamiento & --- \\
\bottomrule
\end{tabular}
\end{mitabla}

Para los empleadores el análisis es solo descriptivo: con cinco
organizaciones, cualquier tabla de contingencia tendría frecuencias
esperadas menores a 5 en todas sus celdas y la prueba chi-cuadrada no
sería confiable.

\subsection{Decisiones de Procesamiento}

El formulario se aplica tal como fue diseñado. Las situaciones que
afectan el cálculo se resuelven con las siguientes decisiones, que la
herramienta aplica de forma automática:

\begin{itemize}
  \item \textbf{Tiempo hasta el primer empleo.} Solo lo responde la rama
  ``trabaja en una organización''. Para calcular su media se reemplaza
  cada rango por su punto medio en meses: ya trabajaba $=0$, menos de un
  mes $=0.5$, 1--4 meses $=2.5$, 4--8 meses $=6$, 8--12 meses $=10$. La
  opción ``más de 12 meses'' no tiene límite superior y se toma como 12,
  por lo que la media obtenida es una cota inferior.
  \item \textbf{Escalas de acuerdo de titulados} (4 niveles). Se
  codifican de 1 (totalmente en desacuerdo) a 4 (totalmente de
  acuerdo). La pregunta de satisfacción usa esta misma escala, por lo
  que se interpreta como grado de acuerdo con estar satisfecho. Para las
  tablas de contingencia los cuatro niveles se agrupan en ``en
  desacuerdo'' y ``de acuerdo'', de modo que las celdas tengan
  suficientes datos.
  \item \textbf{Escala de empleadores.} Tiene dos niveles de acuerdo, uno
  de desacuerdo y la opción ``no sabe''; al no ser simétrica, no se
  calcula media: se reporta el porcentaje de cada opción.
  \item \textbf{Otras agrupaciones para los cruces:} remuneración (hasta
  Bs 5.000 / más de Bs 5.000), tiempo al primer empleo (hasta 4 meses /
  más de 4 meses), institución de posgrado (UMSS: EUPG o FCyT / otra) y
  nivel cursado (diplomado o especialidad / maestría o superior).
  \item \textbf{Cargos.} Las listas de cargos de ambas encuestas son
  distintas; para compararlas se usa una tabla de homologación propia
  (Anexo~\ref{anexo:homologacion}).
  \item \textbf{Datos personales.} Correo, teléfono y LinkedIn se eliminan
  al importar los datos, antes de cualquier cálculo. Las organizaciones
  se listan, pero ningún resultado se atribuye a una organización con
  nombre.
\end{itemize}

\subsection{Pipeline de Procesamiento}

\begin{enumerate}
  \item \textbf{Importación y depuración:} el CSV exportado de cada
  formulario se importa en la herramienta, que elimina los datos
  personales, toma la fecha de cada respuesta y valida que estén las
  columnas necesarias.
  \item \textbf{Selección del período:} se eligen las respuestas
  recibidas entre dos fechas (o todas).
  \item \textbf{Estadística univariada:} frecuencias absolutas y
  relativas de cada pregunta; media, desviación estándar e intervalo de
  confianza del 95\% para las variables numéricas y las escalas.
  \item \textbf{Estadística bivariada:} tablas de contingencia y prueba
  chi-cuadrada para los cruces C1--C7.
  \item \textbf{Simulación:} generación de cohortes de titulados con la
  Transformada Inversa discreta y comparación de escenarios.
  \item \textbf{Visualización y exportación:} panel de control con
  gráficos de barras, mapas de calor y diagrama de radar; exportación de
  los datos limpios y de cada matriz cruzada en CSV.
\end{enumerate}

\subsection{Formulación del Modelo Estadístico}

\subsubsection{Intervalo de Confianza para una Media}\label{sec:ic}

La encuesta no llega a todos los titulados, sino a una muestra. La media
calculada con esa muestra no es exactamente la media de todos los
titulados; el intervalo de confianza indica entre qué valores se
encuentra esta última con una probabilidad del 95\%, y así evita
presentar un valor muestral como si fuera exacto.

Para las variables numéricas (años de vida profesional, años desempleado,
tiempo al primer empleo) y las escalas de acuerdo se
calcula un intervalo de confianza del 95\% para la media poblacional
$\mu$, siguiendo el procedimiento descrito en Coss Bu (1998) para el
caso en que la varianza de la población es desconocida:

\begin{enumerate}
  \item Calcular la media muestral $\bar{x}$, sumando todos los valores
  recolectados y dividiendo entre el tamaño de la muestra $n$.
  \item Calcular la desviación estándar muestral $S$, que mide qué tan
  dispersos están los datos respecto a $\bar{x}$.
  \item Determinar el valor $t_{n-1,\alpha/2}$ de la distribución $t$ de
  Student, buscándolo en tabla según los grados de libertad ($n-1$) y
  el nivel de confianza elegido (95\%, es decir $\alpha=0.05$).
  \item Calcular el margen de error: $\dfrac{S}{\sqrt{n}}\,t_{n-1,\alpha/2}$.
  \item Construir el intervalo sumando y restando el margen de error a
  la media muestral:
\end{enumerate}

\begin{equation}
\bar{x} - \frac{S}{\sqrt{n}}\,t_{n-1,\alpha/2} \;\le\; \mu \;\le\; \bar{x} + \frac{S}{\sqrt{n}}\,t_{n-1,\alpha/2}
\end{equation}

\subsubsection{Prueba de Independencia Chi-Cuadrada}\label{sec:chi}

Al cruzar dos preguntas de la encuesta, las diferencias que aparecen en
la tabla pueden deberse al azar de la muestra. La prueba chi-cuadrada
permite decidir si esas diferencias son lo bastante grandes como para
afirmar que las dos variables están relacionadas, y así no sacar
conclusiones de diferencias casuales.

La prueba de las frecuencias (Chi-Cuadrado) vista en la asignatura
compara, para una variable con $k$ categorías, la frecuencia observada
$FO_i$ contra la frecuencia esperada $FE_i = n/k$ en cada categoría,
mediante el estadístico $X_0^2 = \sum_{i=1}^{k} (FO_i-FE_i)^2/FE_i$,
comparado contra el valor teórico $X_{\alpha,k-1}^2$ de tabla. La
asignatura también cubre la prueba de las series, que extiende esta
misma lógica a pares ordenados dispuestos en una grilla de $k\times k$
celdas ($FO_{ij}$), con frecuencia esperada uniforme
$FE = (n-1)/k^2$ en cada celda.

Para este informe se necesita evaluar si dos variables categóricas
están relacionadas entre sí. El procedimiento se aplica a los cruces
C1--C7 de la Tabla~\ref{tab:cruces}, usando la misma estructura de
grilla $FO_{ij}$ de la prueba de las series. La diferencia es que, a
diferencia de un par de números aleatorios, las categorías de dos
preguntas de la encuesta no tienen por qué estar distribuidas de forma
uniforme entre sí, por lo que la frecuencia esperada de cada celda no
puede asumirse uniforme ($FE=(n-1)/k^2$): debe calcularse a partir de
los totales de fila y columna realmente observados. Esto sigue los
mismos pasos del procedimiento de la asignatura:

\begin{enumerate}
  \item Contar con la muestra: construir la tabla de contingencia,
  cruzando ambas variables. Cada celda $FO_{ij}$ es la cantidad de
  encuestados que respondió a la vez la categoría $i$ de la primera
  variable y la categoría $j$ de la segunda (frecuencia observada), y
  $N$ es el total de encuestados.
  \item Calcular la frecuencia esperada de cada celda, es decir, la
  cantidad que se esperaría si ambas variables fueran independientes
  (generalización de $FE=n/k$ para totales de fila y columna no
  uniformes):
\end{enumerate}

\begin{equation}
FE_{ij} = \frac{(\text{total de la fila } i)(\text{total de la columna } j)}{N}
\end{equation}

\begin{enumerate}
  \setcounter{enumi}{2}
  \item Calcular el estadístico Chi-Cuadrado, sumando la diferencia al
  cuadrado entre lo observado y lo esperado de cada celda, dividida
  entre lo esperado:
\end{enumerate}

\begin{equation}
X_0^2 = \sum_{i=1}^{k}\sum_{j=1}^{m} \frac{(FO_{ij} - FE_{ij})^2}{FE_{ij}}
\end{equation}

\begin{enumerate}
  \setcounter{enumi}{3}
  \item Determinar el estadístico teórico $X_{\alpha,(k-1)(m-1)}^2$ de
  tabla ($k$ y $m$ son el número de categorías de cada variable). Los
  grados de libertad $(k-1)(m-1)$ son mayores que los $k^2-1$ de la
  prueba de las series porque, al estimar $FE_{ij}$ con los totales de
  fila y columna (en vez de asumirla uniforme), se fijan $k-1$ totales
  de fila y $m-1$ totales de columna adicionales.
  \item Establecer la relación entre estadísticos: $X_0^2 <
  X_{\alpha,(k-1)(m-1)}^2$.
  \item Conclusión: si se cumple la relación anterior, no se puede
  rechazar que las variables son independientes; si no se cumple, hay
  evidencia de que están relacionadas.
\end{enumerate}

En resumen, para nuestro caso cada símbolo significa:

\begin{itemize}
  \item $N$: número total de encuestados que entraron a la tabla (en el
  Paso 1).
  \item $k$: número de categorías de la primera variable (ej.
  financiamiento tiene $k=3$: recursos propios, beca total, beca
  parcial).
  \item $m$: número de categorías de la segunda variable (ej. tipo de
  posgrado).
  \item $FO_{ij}$: cantidad real de encuestados que cayó en la
  categoría $i$ de la primera variable \textbf{y} la categoría $j$ de
  la segunda, al mismo tiempo (lo que efectivamente respondieron).
  \item $FE_{ij}$: cantidad que se esperaría en esa misma celda
  \textbf{si no hubiera relación} entre las dos variables (se calcula,
  no se cuenta directamente de la encuesta).
  \item $X_0^2$: un solo número que resume qué tan lejos está lo
  observado de lo esperado en toda la tabla; mientras más grande, más
  evidencia de que las variables sí están relacionadas.
\end{itemize}

El desarrollo completo de esta prueba, paso a paso y con los datos de la
encuesta, se presenta para cada cruce en la Sección~\ref{sec:cruces}.

\subsection{Modelo de Simulación}\label{sec:modelo-sim}

Los dos indicadores literales del criterio 3.2.3 (tasa de empleo y tiempo
medio para obtener el primer empleo) se estudian también por simulación.
El modelo genera cohortes de titulados: a cada titulado simulado se le
asigna una situación laboral y, si trabaja en una organización, un
tiempo hasta su primer empleo. Ambas variables son discretas y su
distribución es la observada en la encuesta, por lo que se generan con
el método de la Transformada Inversa discreta.

\subsubsection{Generación de números aleatorios}

Se usa un generador congruencial mixto:

\begin{equation}
x_{n+1} = (a\,x_n + c) \bmod m, \qquad R = \frac{x_{n+1}}{m}, \qquad 0 \le R < 1
\end{equation}

con $a = 1\,664\,525$, $c = 1\,013\,904\,223$ y $m = 2^{32}$. Se verifican
las tres propiedades de periodo completo:

\begin{enumerate}[label=\roman*)]
  \item $c$ y $m$ son primos entre sí: $m=2^{32}$ solo tiene el factor
  2 y $c$ es impar.
  \item Todo primo que divide a $m$ divide a $a-1$: el único primo es 2
  y $a-1 = 1\,664\,524$ es par.
  \item $m$ es múltiplo de 4 y $a-1 = 4 \times 416\,131$ también lo es.
\end{enumerate}

Como se cumplen i), ii) y iii), el generador tiene periodo completo
($2^{32}$ valores), mucho mayor que la cantidad de números que usa la
simulación.

% Generado por sistema/backend ReporteTex.java -- no editar a mano.
% Datos: dataset de prueba (respuestas sinteticas)
Con $a = 1\,664\,525$, $c = 1\,013\,904\,223$, $m = 2^{32} = 4\,294\,967\,296$ y semilla $x_0 = 12\,345$, las primeras iteraciones son:

\begin{mitabla}{Primeras iteraciones del generador congruencial mixto}
\begin{tabular}{crrrc}
\toprule
$n$ & $x_n$ & $a \cdot x_n + c$ & $x_{n+1} = (a\,x_n + c) \bmod m$ & $R = x_{n+1}/m$ \\
\midrule
0 & 12\,345 & 21\,562\,465\,348 & 87\,628\,868 & 0.0204 \\
1 & 87\,628\,868 & 145\,861\,455\,411\,923 & 71\,072\,467 & 0.0165 \\
2 & 71\,072\,467 & 118\,302\,912\,037\,398 & 2\,332\,836\,374 & 0.5432 \\
3 & 2\,332\,836\,374 & 3\,883\,065\,479\,336\,573 & 2\,726\,892\,157 & 0.6349 \\
4 & 2\,726\,892\,157 & 4\,538\,981\,181\,534\,648 & 3\,908\,547\,000 & 0.9100 \\
5 & 3\,908\,547\,000 & 6\,505\,875\,209\,079\,223 & 483\,019\,191 & 0.1125 \\
\bottomrule
\end{tabular}
\end{mitabla}

Por ejemplo, para $n = 0$: $a \cdot x_0 + c = 1\,664\,525 \times 12\,345 + 1\,013\,904\,223 = 21\,562\,465\,348$; al dividir entre $m$ el residuo es $x_1 = 87\,628\,868$ y $R = 87\,628\,868 / 4\,294\,967\,296 = 0.0204$. Estos son los valores $R_1, R_2, \ldots$ que usa la simulación.



\subsubsection{Transformada Inversa discreta}

\textit{Paso 1. Contar con la función de probabilidad $p(x)$.} Para cada
variable, $p(x)$ es la frecuencia relativa observada de cada categoría
en la encuesta.

\textit{Paso 2. Definir la función acumulada} $F(x_i) = \sum_{j \le i} p(x_j)$.

\textit{Paso 3. Igualar $F(x)$ y $R$.} Cada categoría ocupa el rango de $R$
$[F(x_{i-1}),\,F(x_i))$.

\textit{Paso 4. Aplicar la inversa.} $x = F^{-1}(R)$ es la primera
categoría $x_i$ que cumple $R < F(x_i)$.

Las tablas con las probabilidades obtenidas de la encuesta y sus rangos
de $R$ se presentan en la Sección~\ref{sec:simulacion}.

\textit{Paso 5. Implementar el algoritmo.}

\begin{lstlisting}[caption={Algoritmo de una corrida de la simulación}, language=Pascal]
para cada corrida k = 1 .. n hacer
    empleados <- 0 ; enOrganizacion <- 0 ; sumaMeses <- 0
    para cada titulado t = 1 .. N hacer
        R1 <- gcm() ; R2 <- gcm()
        situacion <- F_situacion^-1(R1)
        si situacion <> "No trabaja" entonces empleados <- empleados + 1
        si situacion = "Trabaja en una organizacion" entonces
            meses <- puntoMedio(F_tiempo^-1(R2))
            enOrganizacion <- enOrganizacion + 1
            sumaMeses <- sumaMeses + meses
        finSi
    finPara
    tasa[k]   <- empleados / N
    tiempo[k] <- sumaMeses / enOrganizacion
finPara
\end{lstlisting}

$R_2$ se genera para todos los titulados aunque no se use, para que el
escenario base y el modificado consuman exactamente la misma secuencia
de números aleatorios y la diferencia entre ambos se deba solo al
cambio de probabilidades.

\subsubsection{Escenarios y número de corridas}

El escenario base usa las probabilidades observadas. El escenario
modificado plantea una situación hipotética que la Carrera podría buscar
con mayor vinculación con empresas: se reduce en un 50\% la probabilidad
de ``no trabaja'' (que pasa a ``trabaja en una organización'') y se
traslada el 30\% de la probabilidad de los tiempos mayores a 4 meses a
la categoría 1--4 meses. Ambos porcentajes son parámetros de la
herramienta.

De cada escenario se ejecutan $n = 1000$ corridas con cohortes de
$N = 100$ titulados. Con los $n$ resultados de cada indicador se calcula
su media y su intervalo de confianza del 95\% con la misma fórmula de la
sección anterior. Como el ancho del intervalo es $2S\,t_{n-1,\alpha/2}/\sqrt{n}$,
reducirlo a la mitad exige cuadruplicar las corridas (Coss Bu, 1998); con
$n=1000$ el margen de error de la tasa de empleo queda por debajo de
0.3 puntos porcentuales.

\subsection{Herramienta de Software}

La herramienta tiene dos partes:

\begin{itemize}
  \item \textbf{Backend en Java}, sin dependencias externas (servidor
  HTTP del propio JDK). Lee los CSV, aplica las decisiones de
  procesamiento, calcula indicadores, cruces y simulación, y expone los
  resultados como servicios REST en formato JSON.
  \item \textbf{Frontend web} en Vue, con la librería de componentes
  Naive UI y la librería de gráficos ECharts. Se organiza por módulos
  (titulados, empleadores, análisis y datos): indicadores de
  empleabilidad y posgrado, valoración de la formación con diagrama de
  radar, cruces de variables con mapas de calor, simulación de
  escenarios con sus tablas de Transformada Inversa, e importación y
  exportación de datos. Todas las pantallas pueden filtrarse por el
  período en que se recibieron las respuestas. Las Historias de Usuario del enunciado se
  cubren con estas funcionalidades (Tabla~\ref{tab:hu}).
\end{itemize}

\begin{mitabla}{Pantallas de la herramienta y servicios que las alimentan}
\begin{tabular}{L{4.4cm}L{4.4cm}L{5.8cm}}
\toprule
Pantalla & Servicio REST & Contenido \\
\midrule
Inicio & \texttt{/api/empleabilidad}, \texttt{/api/formacion}, \texttt{/api/demografia} & Indicadores principales y su criterio ARCU-SUR \\
Empleabilidad & \texttt{/api/empleabilidad} & U1--U5, emprendimiento, razones para no trabajar \\
Formación y posgrado & \texttt{/api/formacion}, \texttt{/api/financiamiento} & U8, U9, fuente de financiamiento \\
Valoración de la formación & \texttt{/api/competencias} & U6, U7 y radar del perfil (U11) \\
Cruces de variables & \texttt{/api/cruces} & C1--C7: mapa de calor, $FO$, $FE$, $X_0^2$ y decisión \\
Perfil y demanda (empleadores) & \texttt{/api/empleadores} & Descriptivos, homologación de cargos, preguntas abiertas \\
Simulación & \texttt{/api/simulacion} & Parámetros, tablas TI, tabla de simulación y corridas \\
Importar y exportar & \texttt{/api/importar}, \texttt{/api/exportar} & Carga de CSV y descarga de datos limpios y matrices \\
\bottomrule
\end{tabular}
\end{mitabla}

\section{Análisis de la Demanda Educativa y Perfil de los Titulados}\label{sec:titulados}

\avisoprueba

Todas las tablas de esta sección y de las dos siguientes se generan con
la herramienta a partir de las respuestas importadas, de modo que cada
número puede rastrearse hasta su cálculo. En las preguntas de opción
única se cuenta cuántos titulados eligieron cada opción ($f_i$) y se
divide entre el número de respuestas ($f_i/n$).

\subsection{Caracterización Demográfica}

La edad y el género no responden por sí mismos a un indicador ARCU-SUR.
Se presentan para describir a quienes respondieron la encuesta, de modo
que los resultados de las secciones siguientes se lean sabiendo a qué
grupo de titulados representan.

% Generado por sistema/backend ReporteTex.java -- no editar a mano.
% Datos: dataset de prueba (respuestas sinteticas)
\begin{mitabla}{Distribución por edad ($n=100$) \datosprueba}
\begin{tabular}{lccc}
\toprule
Edad & $f_i$ & $f_i / n$ & \% \\
\midrule
15–18 & 4 & $4/100 = 0.040$ & 4.0 \\
19–22 & 7 & $7/100 = 0.070$ & 7.0 \\
23–26 & 30 & $30/100 = 0.300$ & 30.0 \\
27–30 & 30 & $30/100 = 0.300$ & 30.0 \\
31–34 & 15 & $15/100 = 0.150$ & 15.0 \\
35 o más & 14 & $14/100 = 0.140$ & 14.0 \\
\midrule
Total & 100 & & 100.0 \\
\bottomrule
\end{tabular}
\end{mitabla}

\textit{Lectura:} la respuesta más frecuente es ``23–26'' ($f_i = 30$, 30.0\% de 100).

\begin{mitabla}{Distribución por género ($n=100$) \datosprueba}
\begin{tabular}{lccc}
\toprule
Género & $f_i$ & $f_i / n$ & \% \\
\midrule
Hombre & 49 & $49/100 = 0.490$ & 49.0 \\
Mujer & 50 & $50/100 = 0.500$ & 50.0 \\
Otro & 1 & $1/100 = 0.010$ & 1.0 \\
\midrule
Total & 100 & & 100.0 \\
\bottomrule
\end{tabular}
\end{mitabla}

\textit{Lectura:} la respuesta más frecuente es ``Mujer'' ($f_i = 50$, 50.0\% de 100).



\subsection{Tasa de Empleo}

La tasa de empleo es uno de los indicadores con que ARCU-SUR evalúa si
la formación se refleja en las posibilidades de empleo de los graduados
(criterio 3.2.3). Se calcula como la proporción de titulados que
trabajan, ya sea en una organización o en un emprendimiento propio.

% Generado por sistema/backend ReporteTex.java -- no editar a mano.
% Datos: Base de datos: 148 importado (titulados) / Base de datos: 53 importado (empleadores)
\textit{Paso 1. Contar con la muestra y determinar su tamaño.} Respondieron la pregunta de situación laboral $n = 148$ titulados.

\textit{Paso 2. Contar la frecuencia de cada respuesta.}

\begin{mitabla}{Situación laboral actual ($n=148$)}
\begin{tabular}{lccc}
\toprule
Situación laboral & $f_i$ & $f_i / n$ & \% \\
\midrule
Trabaja en una organización & 108 & $108/148 = 0.730$ & 73.0 \\
Emprendimiento propio & 15 & $15/148 = 0.101$ & 10.1 \\
No trabaja & 25 & $25/148 = 0.169$ & 16.9 \\
\midrule
Total & 148 & & 100.0 \\
\bottomrule
\end{tabular}
\end{mitabla}

\textit{Lectura:} la respuesta más frecuente es ``Trabaja en una organización'' ($f_i = 108$, 73.0\% de 148).

\textit{Paso 3. Calcular la tasa de empleo.} Se consideran empleados quienes trabajan en una organización y quienes tienen un emprendimiento propio:
\[
\text{Tasa de empleo} = \frac{108 + 15}{148} = \frac{123}{148} = 0.8311 = 83.1\%
\]



\subsection{Tiempo Medio para Obtener el Primer Empleo}

El criterio 3.2.3 pide expresamente el tiempo medio para obtener el
primer empleo. Como la encuesta obtiene ese dato de una muestra y no de
todos los titulados, además de la media se calcula su intervalo de
confianza del 95\%, con el procedimiento de la Sección~\ref{sec:ic}, para
indicar entre qué valores se encuentra la media de toda la población.

% Generado por sistema/backend ReporteTex.java -- no editar a mano.
% Datos: Base de datos: 148 importado (titulados) / Base de datos: 53 importado (empleadores)
\textit{Paso 1. Contar con la muestra y determinar su tamaño.} Solo responden quienes trabajan en una organización. Cada rango se reemplaza por su punto medio en meses ($x_i$); ``más de 12 meses'' se toma como 12, por lo que la media es una cota inferior. El tamaño de la muestra es $n = 101$.

\begin{mitabla}{Cálculo de la media y la desviación estándar de tiempo hasta el primer empleo}
\begin{tabular}{lccccc}
\toprule
Respuesta & $x_i$ & $f_i$ & $f_i\,x_i$ & $(x_i-\bar{x})^2$ & $f_i\,(x_i-\bar{x})^2$ \\
\midrule
Ya trabajaba & 0 & 53 & 0.00 & 6.9362 & 367.6177 \\
$<$ 1 mes & 0.5 & 10 & 5.00 & 4.5525 & 45.5252 \\
1–4 meses & 2.5 & 14 & 35.00 & 0.0179 & 0.2501 \\
4–8 meses & 6 & 6 & 36.00 & 11.3322 & 67.9933 \\
8–12 meses & 10 & 13 & 130.00 & 54.2629 & 705.4179 \\
$>$ 12 meses & 12 & 5 & 60.00 & 87.7283 & 438.6413 \\
\midrule
$\Sigma$ & & 101 & 266.00 & & 1625.4455 \\
\bottomrule
\end{tabular}
\end{mitabla}

\textit{Paso 2. Calcular la media muestral.}
\[
\bar{x} = \frac{\sum f_i\,x_i}{n} = \frac{266.00}{101} = 2.6337\ \text{meses}
\]

\textit{Paso 3. Calcular la desviación estándar muestral} (Coss Bu, ecuación 6.5).
\[
S = \sqrt{\frac{\sum f_i\,(x_i-\bar{x})^2}{n-1}} = \sqrt{\frac{1625.4455}{100}} = 4.0317
\]

\textit{Paso 4. Determinar el estadístico teórico.} Con $\alpha = 0.05$ y $n-1 = 100$ grados de libertad, como los grados de libertad son más de 30 la distribución $t$ se aproxima con la normal, igual que en la prueba de los promedios: $t_{100,\,0.025} \approx z_{0.025} = 1.96$.

\textit{Paso 5. Calcular el margen de error.}
\[
\frac{S}{\sqrt{n}}\,t_{n-1,\,\alpha/2} = \frac{4.0317}{\sqrt{101}}\times 1.960 = 0.4012 \times 1.960 = 0.7863
\]

\textit{Paso 6. Construir el intervalo y concluir.}
\[
2.6337 - 0.7863 \;\le\; \mu \;\le\; 2.6337 + 0.7863 \quad\Longrightarrow\quad 1.85 \le \mu \le 3.42
\]

Con un 95\% de confianza, la media de tiempo hasta el primer empleo de todos los titulados está entre 1.85 y 3.42 meses.



\subsection{Años de Vida Profesional y de Desempleo}

Los años de vida profesional describen la experiencia acumulada de los
titulados que respondieron, y los años de desempleo complementan la
tasa de empleo: una tasa alta puede ocultar períodos largos sin trabajo.
Ambos se estiman con su intervalo de confianza.

\paragraph{Años de vida profesional.}

% Generado por sistema/backend ReporteTex.java -- no editar a mano.
% Datos: dataset de prueba (respuestas sinteticas)
\textit{Paso 1. Contar con la muestra y determinar su tamaño.} Los valores son numéricos y se agrupan por valor repetido ($x_i$ con su frecuencia $f_i$). Se descartan las filas en que los años desempleado superan los años de vida profesional (0 en esta muestra). El tamaño de la muestra es $n = 100$.

\begin{mitabla}{Cálculo de la media y la desviación estándar de años de vida profesional \datosprueba}
\begin{tabular}{lccccc}
\toprule
Respuesta & $x_i$ & $f_i$ & $f_i\,x_i$ & $(x_i-\bar{x})^2$ & $f_i\,(x_i-\bar{x})^2$ \\
\midrule
0 años & 0 & 9 & 0.00 & 32.0356 & 288.3204 \\
1 años & 1 & 13 & 13.00 & 21.7156 & 282.3028 \\
2 años & 2 & 2 & 4.00 & 13.3956 & 26.7912 \\
3 años & 3 & 15 & 45.00 & 7.0756 & 106.1340 \\
4 años & 4 & 15 & 60.00 & 2.7556 & 41.3340 \\
5 años & 5 & 3 & 15.00 & 0.4356 & 1.3068 \\
6 años & 6 & 8 & 48.00 & 0.1156 & 0.9248 \\
7 años & 7 & 7 & 49.00 & 1.7956 & 12.5692 \\
8 años & 8 & 7 & 56.00 & 5.4756 & 38.3292 \\
9 años & 9 & 2 & 18.00 & 11.1556 & 22.3112 \\
10 años & 10 & 3 & 30.00 & 18.8356 & 56.5068 \\
11 años & 11 & 4 & 44.00 & 28.5156 & 114.0624 \\
12 años & 12 & 4 & 48.00 & 40.1956 & 160.7824 \\
13 años & 13 & 2 & 26.00 & 53.8756 & 107.7512 \\
14 años & 14 & 1 & 14.00 & 69.5556 & 69.5556 \\
15 años & 15 & 2 & 30.00 & 87.2356 & 174.4712 \\
20 años & 20 & 1 & 20.00 & 205.6356 & 205.6356 \\
22 años & 22 & 1 & 22.00 & 266.9956 & 266.9956 \\
24 años & 24 & 1 & 24.00 & 336.3556 & 336.3556 \\
\midrule
$\Sigma$ & & 100 & 566.00 & & 2312.4400 \\
\bottomrule
\end{tabular}
\end{mitabla}

\textit{Paso 2. Calcular la media muestral.}
\[
\bar{x} = \frac{\sum f_i\,x_i}{n} = \frac{566.00}{100} = 5.6600\ \text{años}
\]

\textit{Paso 3. Calcular la desviación estándar muestral} (Coss Bu, ecuación 6.5).
\[
S = \sqrt{\frac{\sum f_i\,(x_i-\bar{x})^2}{n-1}} = \sqrt{\frac{2312.4400}{99}} = 4.8330
\]

\textit{Paso 4. Determinar el estadístico teórico.} Con $\alpha = 0.05$ y $n-1 = 99$ grados de libertad, como los grados de libertad son más de 30 la distribución $t$ se aproxima con la normal, igual que en la prueba de los promedios: $t_{99,\,0.025} \approx z_{0.025} = 1.96$.

\textit{Paso 5. Calcular el margen de error.}
\[
\frac{S}{\sqrt{n}}\,t_{n-1,\,\alpha/2} = \frac{4.8330}{\sqrt{100}}\times 1.960 = 0.4833 \times 1.960 = 0.9473
\]

\textit{Paso 6. Construir el intervalo y concluir.}
\[
5.6600 - 0.9473 \;\le\; \mu \;\le\; 5.6600 + 0.9473 \quad\Longrightarrow\quad 4.71 \le \mu \le 6.61
\]

Con un 95\% de confianza, la media de años de vida profesional de todos los titulados está entre 4.71 y 6.61 años. \datosprueba



\paragraph{Años desempleado.}

% Generado por sistema/backend ReporteTex.java -- no editar a mano.
% Datos: dataset de prueba (respuestas sinteticas)
\textit{Paso 1. Contar con la muestra y determinar su tamaño.} Los valores son numéricos y se agrupan por valor repetido ($x_i$ con su frecuencia $f_i$). Se descartan las filas en que los años desempleado superan los años de vida profesional (0 en esta muestra). El tamaño de la muestra es $n = 100$.

\begin{mitabla}{Cálculo de la media y la desviación estándar de años desempleado \datosprueba}
\begin{tabular}{lccccc}
\toprule
Respuesta & $x_i$ & $f_i$ & $f_i\,x_i$ & $(x_i-\bar{x})^2$ & $f_i\,(x_i-\bar{x})^2$ \\
\midrule
0 años & 0 & 42 & 0.00 & 1.0404 & 43.6968 \\
1 años & 1 & 26 & 26.00 & 0.0004 & 0.0104 \\
2 años & 2 & 20 & 40.00 & 0.9604 & 19.2080 \\
3 años & 3 & 12 & 36.00 & 3.9204 & 47.0448 \\
\midrule
$\Sigma$ & & 100 & 102.00 & & 109.9600 \\
\bottomrule
\end{tabular}
\end{mitabla}

\textit{Paso 2. Calcular la media muestral.}
\[
\bar{x} = \frac{\sum f_i\,x_i}{n} = \frac{102.00}{100} = 1.0200\ \text{años}
\]

\textit{Paso 3. Calcular la desviación estándar muestral} (Coss Bu, ecuación 6.5).
\[
S = \sqrt{\frac{\sum f_i\,(x_i-\bar{x})^2}{n-1}} = \sqrt{\frac{109.9600}{99}} = 1.0539
\]

\textit{Paso 4. Determinar el estadístico teórico.} Con $\alpha = 0.05$ y $n-1 = 99$ grados de libertad, como los grados de libertad son más de 30 la distribución $t$ se aproxima con la normal, igual que en la prueba de los promedios: $t_{99,\,0.025} \approx z_{0.025} = 1.96$.

\textit{Paso 5. Calcular el margen de error.}
\[
\frac{S}{\sqrt{n}}\,t_{n-1,\,\alpha/2} = \frac{1.0539}{\sqrt{100}}\times 1.960 = 0.1054 \times 1.960 = 0.2066
\]

\textit{Paso 6. Construir el intervalo y concluir.}
\[
1.0200 - 0.2066 \;\le\; \mu \;\le\; 1.0200 + 0.2066 \quad\Longrightarrow\quad 0.81 \le \mu \le 1.23
\]

Con un 95\% de confianza, la media de años desempleado de todos los titulados está entre 0.81 y 1.23 años. \datosprueba



\subsection{Destino Laboral}

El criterio 3.2.3 incluye el destino laboral y profesional de los
graduados. Conocer en qué sector, rubro y cargo trabajan, y cuánto
perciben, permite a la Carrera contrastar el perfil de egreso con los
puestos que efectivamente ocupan. Solo responden quienes trabajan en una
organización.

% Generado por sistema/backend ReporteTex.java -- no editar a mano.
% Datos: Base de datos: 148 importado (titulados) / Base de datos: 53 importado (empleadores)
\begin{mitabla}{Sector de la organización (trabajan en una organización, $n=108$)}
\begin{tabular}{lccc}
\toprule
Sector & $f_i$ & $f_i / n$ & \% \\
\midrule
Público & 21 & $21/108 = 0.194$ & 19.4 \\
Privado & 87 & $87/108 = 0.806$ & 80.6 \\
\midrule
Total & 108 & & 100.0 \\
\bottomrule
\end{tabular}
\end{mitabla}

\textit{Lectura:} la respuesta más frecuente es ``Privado'' ($f_i = 87$, 80.6\% de 108).

\begin{mitabla}{Rubro de la organización ($n=108$)}
\begin{tabular}{lccc}
\toprule
Rubro & $f_i$ & $f_i / n$ & \% \\
\midrule
Telecomunicaciones & 12 & $12/108 = 0.111$ & 11.1 \\
Educación & 7 & $7/108 = 0.065$ & 6.5 \\
Desarrollo de Software & 67 & $67/108 = 0.620$ & 62.0 \\
Redes de Computadoras & 9 & $9/108 = 0.083$ & 8.3 \\
Salud & 4 & $4/108 = 0.037$ & 3.7 \\
Banca & 9 & $9/108 = 0.083$ & 8.3 \\
Otro & 0 & $0/108 = 0.000$ & 0.0 \\
\midrule
Total & 108 & & 100.0 \\
\bottomrule
\end{tabular}
\end{mitabla}

\textit{Lectura:} la respuesta más frecuente es ``Desarrollo de Software'' ($f_i = 67$, 62.0\% de 108).

\begin{mitabla}{Cargo que desempeña ($n=108$)}
\begin{tabular}{lccc}
\toprule
Cargo & $f_i$ & $f_i / n$ & \% \\
\midrule
Ejecutivo & 17 & $17/108 = 0.157$ & 15.7 \\
Técnico & 13 & $13/108 = 0.120$ & 12.0 \\
Desarrollo & 70 & $70/108 = 0.648$ & 64.8 \\
Help Desk & 8 & $8/108 = 0.074$ & 7.4 \\
Otro & 0 & $0/108 = 0.000$ & 0.0 \\
\midrule
Total & 108 & & 100.0 \\
\bottomrule
\end{tabular}
\end{mitabla}

\textit{Lectura:} la respuesta más frecuente es ``Desarrollo'' ($f_i = 70$, 64.8\% de 108).

\begin{mitabla}{Remuneración mensual ($n=108$)}
\begin{tabular}{lccc}
\toprule
Remuneración & $f_i$ & $f_i / n$ & \% \\
\midrule
$<$ Bs 2.500 & 5 & $5/108 = 0.046$ & 4.6 \\
Bs 2.500–3.500 & 9 & $9/108 = 0.083$ & 8.3 \\
Bs 3.501–5.000 & 20 & $20/108 = 0.185$ & 18.5 \\
Bs 5.001–7.000 & 19 & $19/108 = 0.176$ & 17.6 \\
Bs 7.001–10.000 & 25 & $25/108 = 0.231$ & 23.1 \\
$>$ Bs 10.000 & 30 & $30/108 = 0.278$ & 27.8 \\
\midrule
Total & 108 & & 100.0 \\
\bottomrule
\end{tabular}
\end{mitabla}

\textit{Lectura:} la respuesta más frecuente es ``$>$ Bs 10.000'' ($f_i = 30$, 27.8\% de 108).



\subsection{Formación Continua y Fuentes de Financiamiento}

Los criterios 2.4.1 y 1.1.5 piden que la institución ofrezca y promueva
programas de actualización y posgrado afines a la carrera. Para que esa
oferta responda a una demanda real se necesita saber qué formación ya
tienen los titulados, qué les interesa estudiar, dónde lo harían y cómo
financiaron lo que ya cursaron.

% Generado por sistema/backend ReporteTex.java -- no editar a mano.
% Datos: Base de datos: 148 importado (titulados) / Base de datos: 53 importado (empleadores)
Proporción con formación complementaria: $98/148 = 0.662$ (66.2\%). Proporción con interés en un posgrado: $123/148 = 0.831$ (83.1\%).

\begin{mitabla}{Área de interés para un posgrado (interesados, $n=123$)}
\begin{tabular}{lccc}
\toprule
Área & $f_i$ & $f_i / n$ & \% \\
\midrule
IA & 54 & $54/123 = 0.439$ & 43.9 \\
Robótica & 5 & $5/123 = 0.041$ & 4.1 \\
Ing. de Software & 29 & $29/123 = 0.236$ & 23.6 \\
Base de Datos & 10 & $10/123 = 0.081$ & 8.1 \\
Redes y Seguridad & 25 & $25/123 = 0.203$ & 20.3 \\
\midrule
Total & 123 & & 100.0 \\
\bottomrule
\end{tabular}
\end{mitabla}

\textit{Lectura:} la respuesta más frecuente es ``IA'' ($f_i = 54$, 43.9\% de 123).

\begin{mitabla}{Institución donde estudiaría el posgrado ($n=123$)}
\begin{tabular}{lccc}
\toprule
Institución & $f_i$ & $f_i / n$ & \% \\
\midrule
EUPG & 13 & $13/123 = 0.106$ & 10.6 \\
Dirección de Posgrado FCyT & 38 & $38/123 = 0.309$ & 30.9 \\
IBNORCA & 3 & $3/123 = 0.024$ & 2.4 \\
Univ. internacionales & 31 & $31/123 = 0.252$ & 25.2 \\
Univ. públicas nacionales & 17 & $17/123 = 0.138$ & 13.8 \\
Univ. privadas nacionales & 10 & $10/123 = 0.081$ & 8.1 \\
Otro & 11 & $11/123 = 0.089$ & 8.9 \\
\midrule
Total & 123 & & 100.0 \\
\bottomrule
\end{tabular}
\end{mitabla}

\textit{Lectura:} la respuesta más frecuente es ``Dirección de Posgrado FCyT'' ($f_i = 38$, 30.9\% de 123).

\begin{mitabla}{Nivel del programa cursado (con formación complementaria, $n=98$)}
\begin{tabular}{lccc}
\toprule
Nivel & $f_i$ & $f_i / n$ & \% \\
\midrule
Diplomado & 66 & $66/98 = 0.673$ & 67.3 \\
Especialidad & 3 & $3/98 = 0.031$ & 3.1 \\
Maestría & 27 & $27/98 = 0.276$ & 27.6 \\
Doctorado & 2 & $2/98 = 0.020$ & 2.0 \\
Pos doctorado & 0 & $0/98 = 0.000$ & 0.0 \\
\midrule
Total & 98 & & 100.0 \\
\bottomrule
\end{tabular}
\end{mitabla}

\textit{Lectura:} la respuesta más frecuente es ``Diplomado'' ($f_i = 66$, 67.3\% de 98).

\begin{mitabla}{Fuente de financiamiento del programa cursado ($n=98$)}
\begin{tabular}{lccc}
\toprule
Financiamiento & $f_i$ & $f_i / n$ & \% \\
\midrule
Recursos propios & 90 & $90/98 = 0.918$ & 91.8 \\
Beca total & 3 & $3/98 = 0.031$ & 3.1 \\
Beca parcial & 5 & $5/98 = 0.051$ & 5.1 \\
\midrule
Total & 98 & & 100.0 \\
\bottomrule
\end{tabular}
\end{mitabla}

\textit{Lectura:} la respuesta más frecuente es ``Recursos propios'' ($f_i = 90$, 91.8\% de 98).



\subsection{Valoración de la Formación Recibida}

El criterio 3.2.1 pide evaluar el resultado del proceso formativo, y el
2.1.7 actualizar el currículo según los desarrollos profesionales. La
opinión de los titulados sobre su formación, sobre los aspectos útiles y
a mejorar y sobre las asignaturas es la evidencia directa para ambos.
Las escalas de acuerdo se codifican de 1 a 4 y se les aplica el mismo
procedimiento del intervalo de confianza.

% Generado por sistema/backend ReporteTex.java -- no editar a mano.
% Datos: dataset de prueba (respuestas sinteticas)
\paragraph{Grado de acuerdo con estar satisfecho con la formación.}

\textit{Paso 1. Contar con la muestra y determinar su tamaño.} Cada respuesta se codifica de 1 (totalmente en desacuerdo) a 4 (totalmente de acuerdo). El tamaño de la muestra es $n = 100$.

\begin{mitabla}{Cálculo de la media y la desviación estándar de la satisfacción con la formación (escala 1--4) \datosprueba}
\begin{tabular}{lccccc}
\toprule
Respuesta & $x_i$ & $f_i$ & $f_i\,x_i$ & $(x_i-\bar{x})^2$ & $f_i\,(x_i-\bar{x})^2$ \\
\midrule
Totalmente en desacuerdo & 1 & 8 & 8.00 & 4.8841 & 39.0728 \\
En desacuerdo & 2 & 9 & 18.00 & 1.4641 & 13.1769 \\
De acuerdo & 3 & 37 & 111.00 & 0.0441 & 1.6317 \\
Totalmente de acuerdo & 4 & 46 & 184.00 & 0.6241 & 28.7086 \\
\midrule
$\Sigma$ & & 100 & 321.00 & & 82.5900 \\
\bottomrule
\end{tabular}
\end{mitabla}

\textit{Paso 2. Calcular la media muestral.}
\[
\bar{x} = \frac{\sum f_i\,x_i}{n} = \frac{321.00}{100} = 3.2100\ \text{puntos}
\]

\textit{Paso 3. Calcular la desviación estándar muestral} (Coss Bu, ecuación 6.5).
\[
S = \sqrt{\frac{\sum f_i\,(x_i-\bar{x})^2}{n-1}} = \sqrt{\frac{82.5900}{99}} = 0.9134
\]

\textit{Paso 4. Determinar el estadístico teórico.} Con $\alpha = 0.05$ y $n-1 = 99$ grados de libertad, como los grados de libertad son más de 30 la distribución $t$ se aproxima con la normal, igual que en la prueba de los promedios: $t_{99,\,0.025} \approx z_{0.025} = 1.96$.

\textit{Paso 5. Calcular el margen de error.}
\[
\frac{S}{\sqrt{n}}\,t_{n-1,\,\alpha/2} = \frac{0.9134}{\sqrt{100}}\times 1.960 = 0.0913 \times 1.960 = 0.1790
\]

\textit{Paso 6. Construir el intervalo y concluir.}
\[
3.2100 - 0.1790 \;\le\; \mu \;\le\; 3.2100 + 0.1790 \quad\Longrightarrow\quad 3.03 \le \mu \le 3.39
\]

Con un 95\% de confianza, la media de la satisfacción con la formación (escala 1--4) de todos los titulados está entre 3.03 y 3.39 puntos. \datosprueba

\paragraph{Concordancia entre la formación y los requerimientos del mercado.}

\textit{Paso 1. Contar con la muestra y determinar su tamaño.} Misma codificación de 1 a 4. El tamaño de la muestra es $n = 100$.

\begin{mitabla}{Cálculo de la media y la desviación estándar de la concordancia formación--mercado (escala 1--4) \datosprueba}
\begin{tabular}{lccccc}
\toprule
Respuesta & $x_i$ & $f_i$ & $f_i\,x_i$ & $(x_i-\bar{x})^2$ & $f_i\,(x_i-\bar{x})^2$ \\
\midrule
Totalmente en desacuerdo & 1 & 10 & 10.00 & 3.3856 & 33.8560 \\
En desacuerdo & 2 & 22 & 44.00 & 0.7056 & 15.5232 \\
De acuerdo & 3 & 42 & 126.00 & 0.0256 & 1.0752 \\
Totalmente de acuerdo & 4 & 26 & 104.00 & 1.3456 & 34.9856 \\
\midrule
$\Sigma$ & & 100 & 284.00 & & 85.4400 \\
\bottomrule
\end{tabular}
\end{mitabla}

\textit{Paso 2. Calcular la media muestral.}
\[
\bar{x} = \frac{\sum f_i\,x_i}{n} = \frac{284.00}{100} = 2.8400\ \text{puntos}
\]

\textit{Paso 3. Calcular la desviación estándar muestral} (Coss Bu, ecuación 6.5).
\[
S = \sqrt{\frac{\sum f_i\,(x_i-\bar{x})^2}{n-1}} = \sqrt{\frac{85.4400}{99}} = 0.9290
\]

\textit{Paso 4. Determinar el estadístico teórico.} Con $\alpha = 0.05$ y $n-1 = 99$ grados de libertad, como los grados de libertad son más de 30 la distribución $t$ se aproxima con la normal, igual que en la prueba de los promedios: $t_{99,\,0.025} \approx z_{0.025} = 1.96$.

\textit{Paso 5. Calcular el margen de error.}
\[
\frac{S}{\sqrt{n}}\,t_{n-1,\,\alpha/2} = \frac{0.9290}{\sqrt{100}}\times 1.960 = 0.0929 \times 1.960 = 0.1821
\]

\textit{Paso 6. Construir el intervalo y concluir.}
\[
2.8400 - 0.1821 \;\le\; \mu \;\le\; 2.8400 + 0.1821 \quad\Longrightarrow\quad 2.66 \le \mu \le 3.02
\]

Con un 95\% de confianza, la media de la concordancia formación--mercado (escala 1--4) de todos los titulados está entre 2.66 y 3.02 puntos. \datosprueba

\paragraph{Pertinencia de las labores del puesto actual.}

\textit{Paso 1. Contar con la muestra y determinar su tamaño.} Solo responden quienes trabajan en una organización; misma codificación de 1 a 4. El tamaño de la muestra es $n = 68$.

\begin{mitabla}{Cálculo de la media y la desviación estándar de la pertinencia del puesto (escala 1--4) \datosprueba}
\begin{tabular}{lccccc}
\toprule
Respuesta & $x_i$ & $f_i$ & $f_i\,x_i$ & $(x_i-\bar{x})^2$ & $f_i\,(x_i-\bar{x})^2$ \\
\midrule
Totalmente en desacuerdo & 1 & 8 & 8.00 & 3.6548 & 29.2388 \\
En desacuerdo & 2 & 15 & 30.00 & 0.8313 & 12.4697 \\
De acuerdo & 3 & 20 & 60.00 & 0.0078 & 0.1557 \\
Totalmente de acuerdo & 4 & 25 & 100.00 & 1.1843 & 29.6064 \\
\midrule
$\Sigma$ & & 68 & 198.00 & & 71.4706 \\
\bottomrule
\end{tabular}
\end{mitabla}

\textit{Paso 2. Calcular la media muestral.}
\[
\bar{x} = \frac{\sum f_i\,x_i}{n} = \frac{198.00}{68} = 2.9118\ \text{puntos}
\]

\textit{Paso 3. Calcular la desviación estándar muestral} (Coss Bu, ecuación 6.5).
\[
S = \sqrt{\frac{\sum f_i\,(x_i-\bar{x})^2}{n-1}} = \sqrt{\frac{71.4706}{67}} = 1.0328
\]

\textit{Paso 4. Determinar el estadístico teórico.} Con $\alpha = 0.05$ y $n-1 = 67$ grados de libertad, como los grados de libertad son más de 30 la distribución $t$ se aproxima con la normal, igual que en la prueba de los promedios: $t_{67,\,0.025} \approx z_{0.025} = 1.96$.

\textit{Paso 5. Calcular el margen de error.}
\[
\frac{S}{\sqrt{n}}\,t_{n-1,\,\alpha/2} = \frac{1.0328}{\sqrt{68}}\times 1.960 = 0.1252 \times 1.960 = 0.2455
\]

\textit{Paso 6. Construir el intervalo y concluir.}
\[
2.9118 - 0.2455 \;\le\; \mu \;\le\; 2.9118 + 0.2455 \quad\Longrightarrow\quad 2.67 \le \mu \le 3.16
\]

Con un 95\% de confianza, la media de la pertinencia del puesto (escala 1--4) de todos los titulados está entre 2.67 y 3.16 puntos. \datosprueba

\paragraph{Aspectos de la Carrera, asignaturas.} Son preguntas de elección múltiple: cada titulado puede marcar varias opciones, por lo que $f_i/n$ se calcula sobre el número de titulados ($n = 100$) y los porcentajes pueden sumar más de 100.

\begin{mitabla}{Aspectos de mayor utilidad ($n=100$) \datosprueba}
\begin{tabular}{lccc}
\toprule
Aspecto & $f_i$ & $f_i / n$ & \% \\
\midrule
Distribución de horarios & 55 & $55/100 = 0.550$ & 55.0 \\
Equipos multidisciplinarios & 55 & $55/100 = 0.550$ & 55.0 \\
Pasantías o prácticas & 49 & $49/100 = 0.490$ & 49.0 \\
Habilidades analíticas & 47 & $47/100 = 0.470$ & 47.0 \\
Aplicación práctica en proyectos reales & 46 & $46/100 = 0.460$ & 46.0 \\
Herramientas de la industria & 43 & $43/100 = 0.430$ & 43.0 \\
\midrule
Menciones & 295 & & 295.0 \\
\bottomrule
\end{tabular}
\end{mitabla}

\textit{Lectura:} la respuesta más frecuente es ``Distribución de horarios'' ($f_i = 55$, 55.0\% de 100).

\begin{mitabla}{Aspectos que pueden mejorarse ($n=100$) \datosprueba}
\begin{tabular}{lccc}
\toprule
Aspecto & $f_i$ & $f_i / n$ & \% \\
\midrule
Distribución de horarios & 52 & $52/100 = 0.520$ & 52.0 \\
Actividades extracurriculares & 49 & $49/100 = 0.490$ & 49.0 \\
Relación con empresas (prácticas) & 48 & $48/100 = 0.480$ & 48.0 \\
Actualizar el plan de estudios & 48 & $48/100 = 0.480$ & 48.0 \\
Integrar teoría y práctica & 46 & $46/100 = 0.460$ & 46.0 \\
Más opciones de especialización & 43 & $43/100 = 0.430$ & 43.0 \\
\midrule
Menciones & 286 & & 286.0 \\
\bottomrule
\end{tabular}
\end{mitabla}

\textit{Lectura:} la respuesta más frecuente es ``Distribución de horarios'' ($f_i = 52$, 52.0\% de 100).

\begin{mitabla}{Asignaturas con ventaja competitiva ($n=100$) \datosprueba}
\begin{tabular}{lccc}
\toprule
Asignatura & $f_i$ & $f_i / n$ & \% \\
\midrule
Base de Datos & 65 & $65/100 = 0.650$ & 65.0 \\
Redes de Computadora & 60 & $60/100 = 0.600$ & 60.0 \\
T.I.S. & 59 & $59/100 = 0.590$ & 59.0 \\
Básicas & 57 & $57/100 = 0.570$ & 57.0 \\
Introducción & 55 & $55/100 = 0.550$ & 55.0 \\
Planif. de Proyectos y Gestión & 55 & $55/100 = 0.550$ & 55.0 \\
Tecnología Redes Avanzadas & 52 & $52/100 = 0.520$ & 52.0 \\
Inteligencia Artificial & 49 & $49/100 = 0.490$ & 49.0 \\
Electivas & 48 & $48/100 = 0.480$ & 48.0 \\
\midrule
Menciones & 500 & & 500.0 \\
\bottomrule
\end{tabular}
\end{mitabla}

\textit{Lectura:} la respuesta más frecuente es ``Base de Datos'' ($f_i = 65$, 65.0\% de 100).

\begin{mitabla}{Asignaturas poco útiles ($n=100$) \datosprueba}
\begin{tabular}{lccc}
\toprule
Asignatura & $f_i$ & $f_i / n$ & \% \\
\midrule
Contabilidad Básica & 58 & $58/100 = 0.580$ & 58.0 \\
Taller Prog. Bajo Nivel & 49 & $49/100 = 0.490$ & 49.0 \\
Básicas & 48 & $48/100 = 0.480$ & 48.0 \\
Aplic. Interactivas TV Digital & 47 & $47/100 = 0.470$ & 47.0 \\
Programación Funcional & 47 & $47/100 = 0.470$ & 47.0 \\
Teoría de Grafos & 47 & $47/100 = 0.470$ & 47.0 \\
Telefonía IP & 45 & $45/100 = 0.450$ & 45.0 \\
Teoría de Autómatas & 43 & $43/100 = 0.430$ & 43.0 \\
Web Semánticas & 40 & $40/100 = 0.400$ & 40.0 \\
Robótica & 39 & $39/100 = 0.390$ & 39.0 \\
Graficación por Computadora & 37 & $37/100 = 0.370$ & 37.0 \\
\midrule
Menciones & 500 & & 500.0 \\
\bottomrule
\end{tabular}
\end{mitabla}

\textit{Lectura:} la respuesta más frecuente es ``Contabilidad Básica'' ($f_i = 58$, 58.0\% de 100).



\section{Análisis de la Oferta Laboral y Requerimientos del Sector Productivo}\label{sec:empleadores}

\avisoprueba

El criterio 2.1.2 exige que el perfil de egreso sea coherente con las
competencias que demandan los agentes sociales relevantes, y el 2.4.2
que la Carrera mantenga relaciones con organizaciones públicas y
privadas. Los empleadores son esos agentes: su opinión es la única
fuente del estudio que mira a los titulados desde fuera de la
universidad. Con pocas organizaciones encuestadas el análisis es solo
descriptivo: se presentan las frecuencias de cada respuesta. En la
escala de perfil profesional no se calcula media porque sus opciones no
son simétricas.

% Generado por sistema/backend ReporteTex.java -- no editar a mano.
% Datos: Base de datos: 148 importado (titulados) / Base de datos: 53 importado (empleadores)
\begin{mitabla}{Tipo de organización ($n=53$)}
\begin{tabular}{lccc}
\toprule
Tipo & $f_i$ & $f_i / n$ & \% \\
\midrule
Público & 13 & $13/53 = 0.245$ & 24.5 \\
Privado & 40 & $40/53 = 0.755$ & 75.5 \\
\midrule
Total & 53 & & 100.0 \\
\bottomrule
\end{tabular}
\end{mitabla}

\textit{Lectura:} la respuesta más frecuente es ``Privado'' ($f_i = 40$, 75.5\% de 53).

\begin{mitabla}{Tamaño de la organización ($n=53$)}
\begin{tabular}{lccc}
\toprule
Tamaño & $f_i$ & $f_i / n$ & \% \\
\midrule
Grande (100+) & 20 & $20/53 = 0.377$ & 37.7 \\
Mediana (31–99) & 7 & $7/53 = 0.132$ & 13.2 \\
Pequeña (11–30) & 9 & $9/53 = 0.170$ & 17.0 \\
Micro (hasta 10) & 17 & $17/53 = 0.321$ & 32.1 \\
\midrule
Total & 53 & & 100.0 \\
\bottomrule
\end{tabular}
\end{mitabla}

\textit{Lectura:} la respuesta más frecuente es ``Grande (100+)'' ($f_i = 20$, 37.7\% de 53).

\begin{mitabla}{Contrató Ingenieros de Sistemas en los últimos 5 años ($n=53$)}
\begin{tabular}{lccc}
\toprule
Respuesta & $f_i$ & $f_i / n$ & \% \\
\midrule
Sí & 42 & $42/53 = 0.792$ & 79.2 \\
No & 11 & $11/53 = 0.208$ & 20.8 \\
\midrule
Total & 53 & & 100.0 \\
\bottomrule
\end{tabular}
\end{mitabla}

\textit{Lectura:} la respuesta más frecuente es ``Sí'' ($f_i = 42$, 79.2\% de 53).

\begin{mitabla}{Posibilidad de contratar en los próximos años ($n=53$)}
\begin{tabular}{lccc}
\toprule
Respuesta & $f_i$ & $f_i / n$ & \% \\
\midrule
Alta probabilidad & 16 & $16/53 = 0.302$ & 30.2 \\
Depende del crecimiento & 27 & $27/53 = 0.509$ & 50.9 \\
No prevista & 10 & $10/53 = 0.189$ & 18.9 \\
No requiere & 0 & $0/53 = 0.000$ & 0.0 \\
\midrule
Total & 53 & & 100.0 \\
\bottomrule
\end{tabular}
\end{mitabla}

\textit{Lectura:} la respuesta más frecuente es ``Depende del crecimiento'' ($f_i = 27$, 50.9\% de 53).

\begin{mitabla}{Nivel de formación que demandan ($n=53$)}
\begin{tabular}{lccc}
\toprule
Nivel & $f_i$ & $f_i / n$ & \% \\
\midrule
Licenciatura & 38 & $38/53 = 0.717$ & 71.7 \\
Diplomado & 4 & $4/53 = 0.075$ & 7.5 \\
Especialidad & 8 & $8/53 = 0.151$ & 15.1 \\
Maestría & 3 & $3/53 = 0.057$ & 5.7 \\
Doctorado & 0 & $0/53 = 0.000$ & 0.0 \\
Post-doctorado & 0 & $0/53 = 0.000$ & 0.0 \\
\midrule
Total & 53 & & 100.0 \\
\bottomrule
\end{tabular}
\end{mitabla}

\textit{Lectura:} la respuesta más frecuente es ``Licenciatura'' ($f_i = 38$, 71.7\% de 53).

\begin{mitabla}{Cargos que desempeñan los titulados (elección múltiple, $n=53$)}
\begin{tabular}{lccc}
\toprule
Cargo & $f_i$ & $f_i / n$ & \% \\
\midrule
Técnico & 37 & $37/53 = 0.698$ & 69.8 \\
Analista & 26 & $26/53 = 0.491$ & 49.1 \\
Especialista & 23 & $23/53 = 0.434$ & 43.4 \\
Apoyo / asistencia & 20 & $20/53 = 0.377$ & 37.7 \\
Supervisor / coordinador & 16 & $16/53 = 0.302$ & 30.2 \\
Jefatura / gerencial & 11 & $11/53 = 0.208$ & 20.8 \\
Consultor & 7 & $7/53 = 0.132$ & 13.2 \\
Otro & 0 & $0/53 = 0.000$ & 0.0 \\
\midrule
Menciones & 140 & & 264.2 \\
\bottomrule
\end{tabular}
\end{mitabla}

\textit{Lectura:} la respuesta más frecuente es ``Técnico'' ($f_i = 37$, 69.8\% de 53).

\begin{mitabla}{Formación y perfil profesional: frecuencia $f_i$ y porcentaje $f_i/n$ por opción ($n=53$)}
\begin{tabular}{L{4.8cm}C{2.1cm}C{2.1cm}C{2.3cm}C{1.5cm}}
\toprule
Criterio & Tot.\ de acuerdo & Parc.\ de acuerdo & Tot.\ en desacuerdo & No sabe \\
\midrule
La Carrera brinda confianza como formadora & 29 (55\%) & 18 (34\%) & 4 (8\%) & 2 (4\%) \\
El título es consistente con las competencias & 17 (32\%) & 29 (55\%) & 6 (11\%) & 1 (2\%) \\
Las autoridades consultan nuestra opinión & 7 (13\%) & 18 (34\%) & 16 (30\%) & 12 (23\%) \\
Conozco el perfil profesional declarado & 17 (32\%) & 28 (53\%) & 2 (4\%) & 6 (11\%) \\
El perfil es coherente con el desempeño & 18 (34\%) & 24 (45\%) & 10 (19\%) & 1 (2\%) \\
Desempeño destacado profesional e integral & 19 (36\%) & 28 (53\%) & 5 (9\%) & 1 (2\%) \\
Consideran nuestra opinión en prácticas y pasantías & 14 (26\%) & 17 (32\%) & 9 (17\%) & 13 (25\%) \\
Evidencian competencias laborales & 19 (36\%) & 27 (51\%) & 4 (8\%) & 3 (6\%) \\
Demuestran valores y ética profesional & 24 (45\%) & 20 (38\%) & 5 (9\%) & 4 (8\%) \\
He participado en procesos de retroalimentación & 12 (23\%) & 16 (30\%) & 10 (19\%) & 15 (28\%) \\
La Carrera mantiene vínculos que generan oportunidades & 8 (15\%) & 25 (47\%) & 9 (17\%) & 11 (21\%) \\
He participado en vinculación con el entorno laboral & 5 (9\%) & 20 (38\%) & 20 (38\%) & 8 (15\%) \\
He participado en difusión de oportunidades laborales & 8 (15\%) & 17 (32\%) & 21 (40\%) & 7 (13\%) \\
He participado en vinculación entre profesionales & 4 (8\%) & 23 (43\%) & 20 (38\%) & 6 (11\%) \\
He sido consultado sobre el desempeño de los titulados & 8 (15\%) & 17 (32\%) & 21 (40\%) & 7 (13\%) \\
Han consultado nuestra opinión sobre la pertinencia & 9 (17\%) & 14 (26\%) & 22 (42\%) & 8 (15\%) \\
Recurrimos a la UMSS cuando requerimos profesionales & 11 (21\%) & 21 (40\%) & 14 (26\%) & 7 (13\%) \\
He sido consultado en la autoevaluación de la Carrera & 5 (9\%) & 14 (26\%) & 25 (47\%) & 9 (17\%) \\
\bottomrule
\end{tabular}
\end{mitabla}



\subsection{Áreas, Herramientas y Habilidades Demandadas}

Las tres preguntas abiertas se agruparon en categorías y se contaron
sus menciones. Las áreas emergentes más mencionadas fueron desarrollo
móvil, internet de las cosas y ciberseguridad; las herramientas,
Python y sus librerías de datos; y las habilidades, ética profesional y
aprendizaje autónomo. \datosprueba

\section{Diagnóstico y Síntesis de Resultados}

\avisoprueba

\subsection{Pruebas de Independencia}\label{sec:cruces}

Los indicadores anteriores describen cada variable por separado; los
cruces permiten saber si dos de ellas están relacionadas, lo que ayuda a
explicar los resultados (por ejemplo, si quienes perciben una formación
alineada al mercado consiguen empleo antes). Cada cruce de la
Tabla~\ref{tab:cruces} se resuelve con los seis pasos de la prueba
chi-cuadrada de la Sección~\ref{sec:chi}.

% Generado por sistema/backend ReporteTex.java -- no editar a mano.
% Datos: Base de datos: 148 importado (titulados) / Base de datos: 53 importado (empleadores)
\subsubsection*{C1. Situación laboral $\times$ satisfacción con la formación}

\textit{Propósito.} Determinar si la valoración de la formación recibida cambia según la situación laboral actual del titulado, como evidencia del resultado del proceso formativo frente a las condiciones de empleo. Respalda el criterio ARCU-SUR 3.2.1 / 3.2.3. Se aplica sobre: todos los titulados, con nivel de significancia $\alpha = 0.05$.

\textit{Paso 1. Contar con la muestra: tabla de contingencia} (frecuencia observada $FO_{ij}$, $N = 147$).

\begin{mitabla}{C1: frecuencia observada $FO_{ij}$}
\begin{tabular}{lcc|c}
\toprule
Situación laboral & En desacuerdo & De acuerdo & Total fila \\
\midrule
Trabaja en una organización & 27 & 81 & 108 \\
Emprendimiento propio & 5 & 9 & 14 \\
No trabaja & 6 & 19 & 25 \\
\midrule
Total columna & 38 & 109 & $N=147$ \\
\bottomrule
\end{tabular}
\end{mitabla}

\textit{Paso 2. Calcular la frecuencia esperada} de cada celda, $FE_{ij} = (\text{total fila } i)(\text{total columna } j)/N$. Por ejemplo, para la celda ``Trabaja en una organización -- En desacuerdo'': $FE_{11} = 108 \times 38 / 147 = 27.9184$.

\begin{mitabla}{C1: frecuencia esperada $FE_{ij}$}
\begin{tabular}{lcc}
\toprule
Situación laboral & En desacuerdo & De acuerdo \\
\midrule
Trabaja en una organización & $108\cdot38/147 = 27.92$ & $108\cdot109/147 = 80.08$ \\
Emprendimiento propio & $14\cdot38/147 = 3.62$ & $14\cdot109/147 = 10.38$ \\
No trabaja & $25\cdot38/147 = 6.46$ & $25\cdot109/147 = 18.54$ \\
\bottomrule
\end{tabular}
\end{mitabla}

\textit{Paso 3. Calcular el estadístico chi-cuadrado,} sumando $(FO_{ij}-FE_{ij})^2/FE_{ij}$ de cada celda.

\begin{mitabla}{C1: cálculo de $X_0^2$ celda por celda}
\begin{tabular}{L{5.4cm}cccc}
\toprule
Celda & $FO_{ij}$ & $FE_{ij}$ & $(FO_{ij}-FE_{ij})^2$ & $(FO_{ij}-FE_{ij})^2/FE_{ij}$ \\
\midrule
Trabaja en una organización -- En desacuerdo & 27 & 27.9184 & 0.8434 & 0.0302 \\
Trabaja en una organización -- De acuerdo & 81 & 80.0816 & 0.8434 & 0.0105 \\
Emprendimiento propio -- En desacuerdo & 5 & 3.6190 & 1.9070 & 0.5269 \\
Emprendimiento propio -- De acuerdo & 9 & 10.3810 & 1.9070 & 0.1837 \\
No trabaja -- En desacuerdo & 6 & 6.4626 & 0.2140 & 0.0331 \\
No trabaja -- De acuerdo & 19 & 18.5374 & 0.2140 & 0.0115 \\
\midrule
\multicolumn{4}{r}{$X_0^2 = \sum$} & 0.7960 \\
\bottomrule
\end{tabular}
\end{mitabla}

\textit{Paso 4. Determinar el estadístico teórico.} Con $k = 3$ filas y $m = 2$ columnas, los grados de libertad son $(k-1)(m-1) = (3-1)(2-1) = 2$. De la tabla de chi-cuadrada: $X_{0.05,\,2}^2 = 5.991$.

\textit{Paso 5. Establecer la relación entre estadísticos:} $X_0^2 < X_{0.05,\,2}^2 \;\Rightarrow\; 0.7960 < 5.991$.

\textit{Paso 6. Conclusión.}

\begin{center}
\begin{tabular}{l}
Si ($X_0^2 < X_{0.05,\,2}^2$)\\
\quad \fbox{Independencia}\\
Sino\\
\quad Relación entre las variables\\
FinSi
\end{tabular}
\end{center}

No se puede rechazar que situación laboral y satisfacción sean independientes.

\textit{Observación:} 1 de 6 celdas tienen $FE < 5$, por lo que la prueba pierde confiabilidad y la conclusión debe tomarse con cautela.


% Generado por sistema/backend ReporteTex.java -- no editar a mano.
% Datos: Base de datos: 148 importado (titulados) / Base de datos: 53 importado (empleadores)
\subsubsection*{C2. Sector $\times$ pertinencia del puesto}

\textit{Propósito.} Determinar si la concordancia entre la titulación y el empleo depende del sector, público o privado, en que trabaja el titulado. Respalda el criterio ARCU-SUR 3.2.3. Se aplica sobre: trabajan en una organización, con nivel de significancia $\alpha = 0.05$.

\textit{Paso 1. Contar con la muestra: tabla de contingencia} (frecuencia observada $FO_{ij}$, $N = 108$).

\begin{mitabla}{C2: frecuencia observada $FO_{ij}$}
\begin{tabular}{lcc|c}
\toprule
Sector & En desacuerdo & De acuerdo & Total fila \\
\midrule
Público & 2 & 19 & 21 \\
Privado & 8 & 79 & 87 \\
\midrule
Total columna & 10 & 98 & $N=108$ \\
\bottomrule
\end{tabular}
\end{mitabla}

\textit{Paso 2. Calcular la frecuencia esperada} de cada celda, $FE_{ij} = (\text{total fila } i)(\text{total columna } j)/N$. Por ejemplo, para la celda ``Público -- En desacuerdo'': $FE_{11} = 21 \times 10 / 108 = 1.9444$.

\begin{mitabla}{C2: frecuencia esperada $FE_{ij}$}
\begin{tabular}{lcc}
\toprule
Sector & En desacuerdo & De acuerdo \\
\midrule
Público & $21\cdot10/108 = 1.94$ & $21\cdot98/108 = 19.06$ \\
Privado & $87\cdot10/108 = 8.06$ & $87\cdot98/108 = 78.94$ \\
\bottomrule
\end{tabular}
\end{mitabla}

\textit{Paso 3. Calcular el estadístico chi-cuadrado,} sumando $(FO_{ij}-FE_{ij})^2/FE_{ij}$ de cada celda.

\begin{mitabla}{C2: cálculo de $X_0^2$ celda por celda}
\begin{tabular}{L{5.4cm}cccc}
\toprule
Celda & $FO_{ij}$ & $FE_{ij}$ & $(FO_{ij}-FE_{ij})^2$ & $(FO_{ij}-FE_{ij})^2/FE_{ij}$ \\
\midrule
Público -- En desacuerdo & 2 & 1.9444 & 0.0031 & 0.0016 \\
Público -- De acuerdo & 19 & 19.0556 & 0.0031 & 0.0002 \\
Privado -- En desacuerdo & 8 & 8.0556 & 0.0031 & 0.0004 \\
Privado -- De acuerdo & 79 & 78.9444 & 0.0031 & 0.0000 \\
\midrule
\multicolumn{4}{r}{$X_0^2 = \sum$} & 0.0022 \\
\bottomrule
\end{tabular}
\end{mitabla}

\textit{Paso 4. Determinar el estadístico teórico.} Con $k = 2$ filas y $m = 2$ columnas, los grados de libertad son $(k-1)(m-1) = (2-1)(2-1) = 1$. De la tabla de chi-cuadrada: $X_{0.05,\,1}^2 = 3.841$.

\textit{Paso 5. Establecer la relación entre estadísticos:} $X_0^2 < X_{0.05,\,1}^2 \;\Rightarrow\; 0.0022 < 3.841$.

\textit{Paso 6. Conclusión.}

\begin{center}
\begin{tabular}{l}
Si ($X_0^2 < X_{0.05,\,1}^2$)\\
\quad \fbox{Independencia}\\
Sino\\
\quad Relación entre las variables\\
FinSi
\end{tabular}
\end{center}

No se puede rechazar que sector y pertinencia del puesto sean independientes.

\textit{Observación:} 1 de 4 celdas tienen $FE < 5$, por lo que la prueba pierde confiabilidad y la conclusión debe tomarse con cautela.


% Generado por sistema/backend ReporteTex.java -- no editar a mano.
% Datos: Base de datos: 148 importado (titulados) / Base de datos: 53 importado (empleadores)
\subsubsection*{C3. Formación complementaria $\times$ remuneración}

\textit{Propósito.} Determinar si la formación complementaria se asocia con el nivel de remuneración, como sustento para la oferta de actualización profesional y posgrado. Respalda el criterio ARCU-SUR 2.4.1 / 3.2.3. Se aplica sobre: trabajan en una organización, con nivel de significancia $\alpha = 0.05$.

\textit{Paso 1. Contar con la muestra: tabla de contingencia} (frecuencia observada $FO_{ij}$, $N = 108$).

\begin{mitabla}{C3: frecuencia observada $FO_{ij}$}
\begin{tabular}{lcc|c}
\toprule
Formación complementaria & Hasta Bs 5.000 & Más de Bs 5.000 & Total fila \\
\midrule
Sí & 21 & 52 & 73 \\
No & 13 & 22 & 35 \\
\midrule
Total columna & 34 & 74 & $N=108$ \\
\bottomrule
\end{tabular}
\end{mitabla}

\textit{Paso 2. Calcular la frecuencia esperada} de cada celda, $FE_{ij} = (\text{total fila } i)(\text{total columna } j)/N$. Por ejemplo, para la celda ``Sí -- Hasta Bs 5.000'': $FE_{11} = 73 \times 34 / 108 = 22.9815$.

\begin{mitabla}{C3: frecuencia esperada $FE_{ij}$}
\begin{tabular}{lcc}
\toprule
Formación complementaria & Hasta Bs 5.000 & Más de Bs 5.000 \\
\midrule
Sí & $73\cdot34/108 = 22.98$ & $73\cdot74/108 = 50.02$ \\
No & $35\cdot34/108 = 11.02$ & $35\cdot74/108 = 23.98$ \\
\bottomrule
\end{tabular}
\end{mitabla}

\textit{Paso 3. Calcular el estadístico chi-cuadrado,} sumando $(FO_{ij}-FE_{ij})^2/FE_{ij}$ de cada celda.

\begin{mitabla}{C3: cálculo de $X_0^2$ celda por celda}
\begin{tabular}{L{5.4cm}cccc}
\toprule
Celda & $FO_{ij}$ & $FE_{ij}$ & $(FO_{ij}-FE_{ij})^2$ & $(FO_{ij}-FE_{ij})^2/FE_{ij}$ \\
\midrule
Sí -- Hasta Bs 5.000 & 21 & 22.9815 & 3.9263 & 0.1708 \\
Sí -- Más de Bs 5.000 & 52 & 50.0185 & 3.9263 & 0.0785 \\
No -- Hasta Bs 5.000 & 13 & 11.0185 & 3.9263 & 0.3563 \\
No -- Más de Bs 5.000 & 22 & 23.9815 & 3.9263 & 0.1637 \\
\midrule
\multicolumn{4}{r}{$X_0^2 = \sum$} & 0.7694 \\
\bottomrule
\end{tabular}
\end{mitabla}

\textit{Paso 4. Determinar el estadístico teórico.} Con $k = 2$ filas y $m = 2$ columnas, los grados de libertad son $(k-1)(m-1) = (2-1)(2-1) = 1$. De la tabla de chi-cuadrada: $X_{0.05,\,1}^2 = 3.841$.

\textit{Paso 5. Establecer la relación entre estadísticos:} $X_0^2 < X_{0.05,\,1}^2 \;\Rightarrow\; 0.7694 < 3.841$.

\textit{Paso 6. Conclusión.}

\begin{center}
\begin{tabular}{l}
Si ($X_0^2 < X_{0.05,\,1}^2$)\\
\quad \fbox{Independencia}\\
Sino\\
\quad Relación entre las variables\\
FinSi
\end{tabular}
\end{center}

No se puede rechazar que formación complementaria y remuneración sean independientes.


% Generado por sistema/backend ReporteTex.java -- no editar a mano.
% Datos: dataset de prueba (respuestas sinteticas)
\subsubsection*{C4. Interés en posgrado $\times$ oferta de posgrado FCyT}

\textit{Propósito.} Determinar si la percepción de que la oferta de posgrado de la FCyT es afín a la carrera depende del interés en realizar un posgrado. Respalda el criterio ARCU-SUR 2.4.1 / 1.1.5. Se aplica sobre: todos los titulados, con nivel de significancia $\alpha = 0.05$.

\textit{Paso 1. Contar con la muestra: tabla de contingencia} (frecuencia observada $FO_{ij}$, $N = 100$).

\begin{mitabla}{C4: frecuencia observada $FO_{ij}$ \datosprueba}
\begin{tabular}{lcc|c}
\toprule
Interés en posgrado & En desacuerdo & De acuerdo & Total fila \\
\midrule
Sí & 15 & 49 & 64 \\
No & 11 & 25 & 36 \\
\midrule
Total columna & 26 & 74 & $N=100$ \\
\bottomrule
\end{tabular}
\end{mitabla}

\textit{Paso 2. Calcular la frecuencia esperada} de cada celda, $FE_{ij} = (\text{total fila } i)(\text{total columna } j)/N$. Por ejemplo, para la celda ``Sí -- En desacuerdo'': $FE_{11} = 64 \times 26 / 100 = 16.6400$.

\begin{mitabla}{C4: frecuencia esperada $FE_{ij}$ \datosprueba}
\begin{tabular}{lcc}
\toprule
Interés en posgrado & En desacuerdo & De acuerdo \\
\midrule
Sí & $64\cdot26/100 = 16.64$ & $64\cdot74/100 = 47.36$ \\
No & $36\cdot26/100 = 9.36$ & $36\cdot74/100 = 26.64$ \\
\bottomrule
\end{tabular}
\end{mitabla}

\textit{Paso 3. Calcular el estadístico chi-cuadrado,} sumando $(FO_{ij}-FE_{ij})^2/FE_{ij}$ de cada celda.

\begin{mitabla}{C4: cálculo de $X_0^2$ celda por celda \datosprueba}
\begin{tabular}{L{5.4cm}cccc}
\toprule
Celda & $FO_{ij}$ & $FE_{ij}$ & $(FO_{ij}-FE_{ij})^2$ & $(FO_{ij}-FE_{ij})^2/FE_{ij}$ \\
\midrule
Sí -- En desacuerdo & 15 & 16.6400 & 2.6896 & 0.1616 \\
Sí -- De acuerdo & 49 & 47.3600 & 2.6896 & 0.0568 \\
No -- En desacuerdo & 11 & 9.3600 & 2.6896 & 0.2874 \\
No -- De acuerdo & 25 & 26.6400 & 2.6896 & 0.1010 \\
\midrule
\multicolumn{4}{r}{$X_0^2 = \sum$} & 0.6067 \\
\bottomrule
\end{tabular}
\end{mitabla}

\textit{Paso 4. Determinar el estadístico teórico.} Con $k = 2$ filas y $m = 2$ columnas, los grados de libertad son $(k-1)(m-1) = (2-1)(2-1) = 1$. De la tabla de chi-cuadrada: $X_{0.05,\,1}^2 = 3.841$.

\textit{Paso 5. Establecer la relación entre estadísticos:} $X_0^2 < X_{0.05,\,1}^2 \;\Rightarrow\; 0.6067 < 3.841$.

\textit{Paso 6. Conclusión.}

\begin{center}
\begin{tabular}{l}
Si ($X_0^2 < X_{0.05,\,1}^2$)\\
\quad \fbox{Independencia}\\
Sino\\
\quad Relación entre las variables\\
FinSi
\end{tabular}
\end{center}

No se puede rechazar que interés en posgrado y oferta fcyt afín sean independientes. \datosprueba


% Generado por sistema/backend ReporteTex.java -- no editar a mano.
% Datos: Base de datos: 148 importado (titulados) / Base de datos: 53 importado (empleadores)
\subsubsection*{C5. Área de interés $\times$ institución elegida}

\textit{Propósito.} Determinar si la preferencia por estudiar el posgrado en la UMSS depende del área de interés, para identificar áreas en que la oferta propia no capta la demanda. Respalda el criterio ARCU-SUR 2.4.1 / 1.1.5. Se aplica sobre: interesados en posgrado, con nivel de significancia $\alpha = 0.05$.

\textit{Paso 1. Contar con la muestra: tabla de contingencia} (frecuencia observada $FO_{ij}$, $N = 123$).

\begin{mitabla}{C5: frecuencia observada $FO_{ij}$}
\begin{tabular}{lcc|c}
\toprule
Área de interés & UMSS (EUPG o FCyT) & Otra institución & Total fila \\
\midrule
IA & 20 & 34 & 54 \\
Robótica & 2 & 3 & 5 \\
Ing. de Software & 13 & 16 & 29 \\
Base de Datos & 6 & 4 & 10 \\
Redes y Seguridad & 10 & 15 & 25 \\
\midrule
Total columna & 51 & 72 & $N=123$ \\
\bottomrule
\end{tabular}
\end{mitabla}

\textit{Paso 2. Calcular la frecuencia esperada} de cada celda, $FE_{ij} = (\text{total fila } i)(\text{total columna } j)/N$. Por ejemplo, para la celda ``IA -- UMSS (EUPG o FCyT)'': $FE_{11} = 54 \times 51 / 123 = 22.3902$.

\begin{mitabla}{C5: frecuencia esperada $FE_{ij}$}
\begin{tabular}{lcc}
\toprule
Área de interés & UMSS (EUPG o FCyT) & Otra institución \\
\midrule
IA & $54\cdot51/123 = 22.39$ & $54\cdot72/123 = 31.61$ \\
Robótica & $5\cdot51/123 = 2.07$ & $5\cdot72/123 = 2.93$ \\
Ing. de Software & $29\cdot51/123 = 12.02$ & $29\cdot72/123 = 16.98$ \\
Base de Datos & $10\cdot51/123 = 4.15$ & $10\cdot72/123 = 5.85$ \\
Redes y Seguridad & $25\cdot51/123 = 10.37$ & $25\cdot72/123 = 14.63$ \\
\bottomrule
\end{tabular}
\end{mitabla}

\textit{Paso 3. Calcular el estadístico chi-cuadrado,} sumando $(FO_{ij}-FE_{ij})^2/FE_{ij}$ de cada celda.

\begin{mitabla}{C5: cálculo de $X_0^2$ celda por celda}
\begin{tabular}{L{5.4cm}cccc}
\toprule
Celda & $FO_{ij}$ & $FE_{ij}$ & $(FO_{ij}-FE_{ij})^2$ & $(FO_{ij}-FE_{ij})^2/FE_{ij}$ \\
\midrule
IA -- UMSS (EUPG o FCyT) & 20 & 22.3902 & 5.7133 & 0.2552 \\
IA -- Otra institución & 34 & 31.6098 & 5.7133 & 0.1807 \\
Robótica -- UMSS (EUPG o FCyT) & 2 & 2.0732 & 0.0054 & 0.0026 \\
Robótica -- Otra institución & 3 & 2.9268 & 0.0054 & 0.0018 \\
Ing. de Software -- UMSS (EUPG o FCyT) & 13 & 12.0244 & 0.9518 & 0.0792 \\
Ing. de Software -- Otra institución & 16 & 16.9756 & 0.9518 & 0.0561 \\
Base de Datos -- UMSS (EUPG o FCyT) & 6 & 4.1463 & 3.4360 & 0.8287 \\
Base de Datos -- Otra institución & 4 & 5.8537 & 3.4360 & 0.5870 \\
Redes y Seguridad -- UMSS (EUPG o FCyT) & 10 & 10.3659 & 0.1338 & 0.0129 \\
Redes y Seguridad -- Otra institución & 15 & 14.6341 & 0.1338 & 0.0091 \\
\midrule
\multicolumn{4}{r}{$X_0^2 = \sum$} & 2.0133 \\
\bottomrule
\end{tabular}
\end{mitabla}

\textit{Paso 4. Determinar el estadístico teórico.} Con $k = 5$ filas y $m = 2$ columnas, los grados de libertad son $(k-1)(m-1) = (5-1)(2-1) = 4$. De la tabla de chi-cuadrada: $X_{0.05,\,4}^2 = 9.488$.

\textit{Paso 5. Establecer la relación entre estadísticos:} $X_0^2 < X_{0.05,\,4}^2 \;\Rightarrow\; 2.0133 < 9.488$.

\textit{Paso 6. Conclusión.}

\begin{center}
\begin{tabular}{l}
Si ($X_0^2 < X_{0.05,\,4}^2$)\\
\quad \fbox{Independencia}\\
Sino\\
\quad Relación entre las variables\\
FinSi
\end{tabular}
\end{center}

No se puede rechazar que área de interés y institución sean independientes.

\textit{Observación:} 3 de 10 celdas tienen $FE < 5$, por lo que la prueba pierde confiabilidad y la conclusión debe tomarse con cautela.


% Generado por sistema/backend ReporteTex.java -- no editar a mano.
% Datos: dataset de prueba (respuestas sinteticas)
\subsubsection*{C6. Concordancia formación–mercado $\times$ tiempo al primer empleo}

\textit{Propósito.} Determinar si percibir la formación alineada con los requerimientos del mercado se asocia con un menor tiempo para obtener el primer empleo. Respalda el criterio ARCU-SUR 3.2.3. Se aplica sobre: trabajan en una organización, con nivel de significancia $\alpha = 0.05$.

\textit{Paso 1. Contar con la muestra: tabla de contingencia} (frecuencia observada $FO_{ij}$, $N = 68$).

\begin{mitabla}{C6: frecuencia observada $FO_{ij}$ \datosprueba}
\begin{tabular}{lcc|c}
\toprule
Concordancia & Hasta 4 meses & Más de 4 meses & Total fila \\
\midrule
En desacuerdo & 5 & 13 & 18 \\
De acuerdo & 29 & 21 & 50 \\
\midrule
Total columna & 34 & 34 & $N=68$ \\
\bottomrule
\end{tabular}
\end{mitabla}

\textit{Paso 2. Calcular la frecuencia esperada} de cada celda, $FE_{ij} = (\text{total fila } i)(\text{total columna } j)/N$. Por ejemplo, para la celda ``En desacuerdo -- Hasta 4 meses'': $FE_{11} = 18 \times 34 / 68 = 9.0000$.

\begin{mitabla}{C6: frecuencia esperada $FE_{ij}$ \datosprueba}
\begin{tabular}{lcc}
\toprule
Concordancia & Hasta 4 meses & Más de 4 meses \\
\midrule
En desacuerdo & $18\cdot34/68 = 9.00$ & $18\cdot34/68 = 9.00$ \\
De acuerdo & $50\cdot34/68 = 25.00$ & $50\cdot34/68 = 25.00$ \\
\bottomrule
\end{tabular}
\end{mitabla}

\textit{Paso 3. Calcular el estadístico chi-cuadrado,} sumando $(FO_{ij}-FE_{ij})^2/FE_{ij}$ de cada celda.

\begin{mitabla}{C6: cálculo de $X_0^2$ celda por celda \datosprueba}
\begin{tabular}{L{5.4cm}cccc}
\toprule
Celda & $FO_{ij}$ & $FE_{ij}$ & $(FO_{ij}-FE_{ij})^2$ & $(FO_{ij}-FE_{ij})^2/FE_{ij}$ \\
\midrule
En desacuerdo -- Hasta 4 meses & 5 & 9.0000 & 16.0000 & 1.7778 \\
En desacuerdo -- Más de 4 meses & 13 & 9.0000 & 16.0000 & 1.7778 \\
De acuerdo -- Hasta 4 meses & 29 & 25.0000 & 16.0000 & 0.6400 \\
De acuerdo -- Más de 4 meses & 21 & 25.0000 & 16.0000 & 0.6400 \\
\midrule
\multicolumn{4}{r}{$X_0^2 = \sum$} & 4.8356 \\
\bottomrule
\end{tabular}
\end{mitabla}

\textit{Paso 4. Determinar el estadístico teórico.} Con $k = 2$ filas y $m = 2$ columnas, los grados de libertad son $(k-1)(m-1) = (2-1)(2-1) = 1$. De la tabla de chi-cuadrada: $X_{0.05,\,1}^2 = 3.841$.

\textit{Paso 5. Establecer la relación entre estadísticos:} $X_0^2 \ge X_{0.05,\,1}^2 \;\Rightarrow\; 4.8356 \ge 3.841$.

\textit{Paso 6. Conclusión.}

\begin{center}
\begin{tabular}{l}
Si ($X_0^2 < X_{0.05,\,1}^2$)\\
\quad Independencia\\
Sino\\
\quad \fbox{Relación entre las variables}\\
FinSi
\end{tabular}
\end{center}

Se rechaza la independencia: hay evidencia de que concordancia y tiempo al primer empleo están relacionadas. Proporción por fila en la primera columna (``Hasta 4 meses''): En desacuerdo $5/18 = 27.8\%$; De acuerdo $29/50 = 58.0\%$. \datosprueba


% Generado por sistema/backend ReporteTex.java -- no editar a mano.
% Datos: Base de datos: 148 importado (titulados) / Base de datos: 53 importado (empleadores)
\subsubsection*{C7. Financiamiento $\times$ nivel del programa cursado (opcional)}

\textit{Propósito.} Determinar si la fuente de financiamiento depende del nivel del programa cursado; complementa la historia de usuario de financiamiento. Respalda el criterio ARCU-SUR sin indicador ARCU-SUR directo. Se aplica sobre: con formación complementaria, con nivel de significancia $\alpha = 0.05$.

\textit{Paso 1. Contar con la muestra: tabla de contingencia} (frecuencia observada $FO_{ij}$, $N = 98$).

\begin{mitabla}{C7: frecuencia observada $FO_{ij}$}
\begin{tabular}{lcc|c}
\toprule
Financiamiento & Diplomado o Especialidad & Maestría o superior & Total fila \\
\midrule
Recursos propios & 66 & 24 & 90 \\
Beca total & 1 & 2 & 3 \\
Beca parcial & 2 & 3 & 5 \\
\midrule
Total columna & 69 & 29 & $N=98$ \\
\bottomrule
\end{tabular}
\end{mitabla}

\textit{Paso 2. Calcular la frecuencia esperada} de cada celda, $FE_{ij} = (\text{total fila } i)(\text{total columna } j)/N$. Por ejemplo, para la celda ``Recursos propios -- Diplomado o Especialidad'': $FE_{11} = 90 \times 69 / 98 = 63.3673$.

\begin{mitabla}{C7: frecuencia esperada $FE_{ij}$}
\begin{tabular}{lcc}
\toprule
Financiamiento & Diplomado o Especialidad & Maestría o superior \\
\midrule
Recursos propios & $90\cdot69/98 = 63.37$ & $90\cdot29/98 = 26.63$ \\
Beca total & $3\cdot69/98 = 2.11$ & $3\cdot29/98 = 0.89$ \\
Beca parcial & $5\cdot69/98 = 3.52$ & $5\cdot29/98 = 1.48$ \\
\bottomrule
\end{tabular}
\end{mitabla}

\textit{Paso 3. Calcular el estadístico chi-cuadrado,} sumando $(FO_{ij}-FE_{ij})^2/FE_{ij}$ de cada celda.

\begin{mitabla}{C7: cálculo de $X_0^2$ celda por celda}
\begin{tabular}{L{5.4cm}cccc}
\toprule
Celda & $FO_{ij}$ & $FE_{ij}$ & $(FO_{ij}-FE_{ij})^2$ & $(FO_{ij}-FE_{ij})^2/FE_{ij}$ \\
\midrule
Recursos propios -- Diplomado o Especialidad & 66 & 63.3673 & 6.9309 & 0.1094 \\
Recursos propios -- Maestría o superior & 24 & 26.6327 & 6.9309 & 0.2602 \\
Beca total -- Diplomado o Especialidad & 1 & 2.1122 & 1.2371 & 0.5857 \\
Beca total -- Maestría o superior & 2 & 0.8878 & 1.2371 & 1.3935 \\
Beca parcial -- Diplomado o Especialidad & 2 & 3.5204 & 2.3116 & 0.6566 \\
Beca parcial -- Maestría o superior & 3 & 1.4796 & 2.3116 & 1.5624 \\
\midrule
\multicolumn{4}{r}{$X_0^2 = \sum$} & 4.5678 \\
\bottomrule
\end{tabular}
\end{mitabla}

\textit{Paso 4. Determinar el estadístico teórico.} Con $k = 3$ filas y $m = 2$ columnas, los grados de libertad son $(k-1)(m-1) = (3-1)(2-1) = 2$. De la tabla de chi-cuadrada: $X_{0.05,\,2}^2 = 5.991$.

\textit{Paso 5. Establecer la relación entre estadísticos:} $X_0^2 < X_{0.05,\,2}^2 \;\Rightarrow\; 4.5678 < 5.991$.

\textit{Paso 6. Conclusión.}

\begin{center}
\begin{tabular}{l}
Si ($X_0^2 < X_{0.05,\,2}^2$)\\
\quad \fbox{Independencia}\\
Sino\\
\quad Relación entre las variables\\
FinSi
\end{tabular}
\end{center}

No se puede rechazar que financiamiento y nivel cursado sean independientes.

\textit{Observación:} 4 de 6 celdas tienen $FE < 5$, por lo que la prueba pierde confiabilidad y la conclusión debe tomarse con cautela.



\subsubsection*{Síntesis de los cruces}

% Generado por sistema/backend ReporteTex.java -- no editar a mano.
% Datos: Base de datos: 148 importado (titulados) / Base de datos: 53 importado (empleadores)
\begin{mitabla}{Resumen de las pruebas de independencia ($\alpha = 0.05$)}
\begin{tabular}{clcccccc}
\toprule
Id & Población & $N$ & $X_0^2$ & gl & $X_{0.05,\,gl}^2$ & $FE<5$ & Conclusión \\
\midrule
C1 & Todos los titulados & 147 & 0.796 & 2 & 5.991 & 1/6 & Independientes \\
C2 & Trabajan en una organización & 108 & 0.002 & 1 & 3.841 & 1/4 & Independientes \\
C3 & Trabajan en una organización & 108 & 0.769 & 1 & 3.841 & 0/4 & Independientes \\
C4 & Todos los titulados & 148 & 1.749 & 1 & 3.841 & 0/4 & Independientes \\
C5 & Interesados en posgrado & 123 & 2.013 & 4 & 9.488 & 3/10 & Independientes \\
C6 & Trabajan en una organización & 74 & 0.227 & 1 & 3.841 & 0/4 & Independientes \\
C7 & Con formación complementaria & 98 & 4.568 & 2 & 5.991 & 4/6 & Independientes \\
\bottomrule
\end{tabular}
\end{mitabla}



\subsection{Resultados de la Simulación}\label{sec:simulacion}

Los indicadores de la encuesta muestran la situación actual. La
simulación permite estimar cómo cambiarían la tasa de empleo y el
tiempo al primer empleo si la Carrera lograra mejoras concretas, lo que
sirve de sustento cuantitativo para las acciones del plan de mejoras
que pide la autoevaluación. Se aplica el modelo de la
Sección~\ref{sec:modelo-sim} con $n = 1000$
corridas, cohortes del mismo tamaño que la muestra de titulados, semilla
$x_0 = 12345$ y el escenario modificado descrito allí.

\subsubsection*{Tablas de Transformada Inversa}

% Generado por sistema/backend ReporteTex.java -- no editar a mano.
% Datos: dataset de prueba (respuestas sinteticas)
\textit{Paso 1. Contar con la función de probabilidad} $p(x) = f_i/n$, a partir de las frecuencias observadas en la encuesta. \textit{Paso 2. Definir la acumulada} $F(x_i) = F(x_{i-1}) + p(x_i)$. \textit{Pasos 3 y 4. Igualar $F(x) = R$ e invertir:} a cada categoría le corresponde el rango $[F(x_{i-1}),\,F(x_i))$, y $x = F^{-1}(R)$ es la categoría cuyo rango contiene a $R$.

\begin{mitabla}{Transformada Inversa de la situación laboral, escenario base ($n=100$) \datosprueba}
\begin{tabular}{lcccc}
\toprule
$x$ & $f_i$ & $p(x)$ & $F(x)$ & Rango de $R$ \\
\midrule
Trabaja en una organización & 68 & $68/100 = 0.680$ & 0.680 & $[0.000,\ 0.680)$ \\
Emprendimiento propio & 19 & $19/100 = 0.190$ & 0.870 & $[0.680,\ 0.870)$ \\
No trabaja & 13 & $13/100 = 0.130$ & 1.000 & $[0.870,\ 1.000)$ \\
\bottomrule
\end{tabular}
\end{mitabla}

\begin{mitabla}{Transformada Inversa del tiempo al primer empleo, escenario base ($n=68$) \datosprueba}
\begin{tabular}{lccccc}
\toprule
$x$ & $f_i$ & Meses & $p(x)$ & $F(x)$ & Rango de $R$ \\
\midrule
Ya trabajaba & 7 & 0 & $7/68 = 0.103$ & 0.103 & $[0.000,\ 0.103)$ \\
$<$ 1 mes & 12 & 0.5 & $12/68 = 0.176$ & 0.279 & $[0.103,\ 0.279)$ \\
1–4 meses & 15 & 2.5 & $15/68 = 0.221$ & 0.500 & $[0.279,\ 0.500)$ \\
4–8 meses & 11 & 6 & $11/68 = 0.162$ & 0.662 & $[0.500,\ 0.662)$ \\
8–12 meses & 12 & 10 & $12/68 = 0.176$ & 0.838 & $[0.662,\ 0.838)$ \\
$>$ 12 meses & 11 & 12 & $11/68 = 0.162$ & 1.000 & $[0.838,\ 1.000)$ \\
\bottomrule
\end{tabular}
\end{mitabla}

\textit{Escenario modificado.} Se traslada el 50\% de la probabilidad de ``No trabaja'' a ``Trabaja en una organización'':
\[
p^{\prime}(\text{No trabaja}) = 0.130 \times (1 - 0.50) = 0.0650, \qquad p^{\prime}(\text{Trabaja}) = 0.680 + 0.130 \times 0.50 = 0.7450
\]
y el 30\% de la probabilidad de los tiempos mayores a 4 meses a ``1--4 meses'':
\[
p^{\prime}(\text{1--4 meses}) = 0.221 + (0.162 + 0.176 + 0.162) \times 0.30 = 0.3706
\]
y cada uno de esos tres rangos queda multiplicado por $1 - 0.30$.

\begin{mitabla}{Transformada Inversa de la situación laboral, escenario modificado \datosprueba}
\begin{tabular}{lccc}
\toprule
$x$ & $p(x)$ & $F(x)$ & Rango de $R$ \\
\midrule
Trabaja en una organización & 0.745 & 0.745 & $[0.000,\ 0.745)$ \\
Emprendimiento propio & 0.190 & 0.935 & $[0.745,\ 0.935)$ \\
No trabaja & 0.065 & 1.000 & $[0.935,\ 1.000)$ \\
\bottomrule
\end{tabular}
\end{mitabla}

\begin{mitabla}{Transformada Inversa del tiempo al primer empleo, escenario modificado \datosprueba}
\begin{tabular}{lcccc}
\toprule
$x$ & Meses & $p(x)$ & $F(x)$ & Rango de $R$ \\
\midrule
Ya trabajaba & 0 & 0.103 & 0.103 & $[0.000,\ 0.103)$ \\
$<$ 1 mes & 0.5 & 0.176 & 0.279 & $[0.103,\ 0.279)$ \\
1–4 meses & 2.5 & 0.371 & 0.650 & $[0.279,\ 0.650)$ \\
4–8 meses & 6 & 0.113 & 0.763 & $[0.650,\ 0.763)$ \\
8–12 meses & 10 & 0.124 & 0.887 & $[0.763,\ 0.887)$ \\
$>$ 12 meses & 12 & 0.113 & 1.000 & $[0.887,\ 1.000)$ \\
\bottomrule
\end{tabular}
\end{mitabla}



\subsubsection*{Tabla de simulación}

% Generado por sistema/backend ReporteTex.java -- no editar a mano.
% Datos: dataset de prueba (respuestas sinteticas)
\begin{mitabla}{Tabla de simulación: primeros 10 titulados de la corrida 1, escenario base \datosprueba}
\scriptsize
\begin{tabular}{cclcll}
\toprule
$t$ & $R_1 \leftarrow \text{gcm}()$ & Situación $= F^{-1}(R_1)$ & $R_2 \leftarrow \text{gcm}()$ & Tiempo $= F^{-1}(R_2)$ & Meses \\
\midrule
1 & 0.0204 & Trabaja en una organización {\scriptsize ($R_1 < 0.680$)} & 0.0165 & Ya trabajaba {\scriptsize ($R_2 < 0.103$)} & 0 \\
2 & 0.5432 & Trabaja en una organización {\scriptsize ($R_1 < 0.680$)} & 0.6349 & 4–8 meses {\scriptsize ($R_2 < 0.662$)} & 6 \\
3 & 0.9100 & No trabaja {\scriptsize ($R_1 < 1.000$)} & --- & --- & --- \\
4 & 0.4959 & Trabaja en una organización {\scriptsize ($R_1 < 0.680$)} & 0.5483 & 4–8 meses {\scriptsize ($R_2 < 0.662$)} & 6 \\
5 & 0.5961 & Trabaja en una organización {\scriptsize ($R_1 < 0.680$)} & 0.7835 & 8–12 meses {\scriptsize ($R_2 < 0.838$)} & 10 \\
6 & 0.0399 & Trabaja en una organización {\scriptsize ($R_1 < 0.680$)} & 0.7438 & 8–12 meses {\scriptsize ($R_2 < 0.838$)} & 10 \\
7 & 0.9658 & No trabaja {\scriptsize ($R_1 < 1.000$)} & --- & --- & --- \\
8 & 0.6996 & Emprendimiento propio {\scriptsize ($R_1 < 0.870$)} & --- & --- & --- \\
9 & 0.3965 & Trabaja en una organización {\scriptsize ($R_1 < 0.680$)} & 0.9595 & $>$ 12 meses {\scriptsize ($R_2 < 1.000$)} & 12 \\
10 & 0.4061 & Trabaja en una organización {\scriptsize ($R_1 < 0.680$)} & 0.9528 & $>$ 12 meses {\scriptsize ($R_2 < 1.000$)} & 12 \\
\bottomrule
\end{tabular}
\end{mitabla}

Entre paréntesis se indica el límite superior del rango que contiene a $R$: se toma la primera categoría cuya acumulada $F(x)$ es mayor que $R$.



\subsubsection*{Corridas y comparación de escenarios}

% Generado por sistema/backend ReporteTex.java -- no editar a mano.
% Datos: dataset de prueba (respuestas sinteticas)
En cada corrida se simulan $N = 100$ titulados y se calcula:
\[
\text{Tasa de empleo}_k = \frac{\text{empleados}_k}{N}, \qquad \text{Tiempo medio}_k = \frac{\sum \text{meses}}{\text{titulados en una organización}_k}
\]

\begin{mitabla}{Prueba de escritorio: primeras 25 corridas del escenario base ($N=100$) \datosprueba}
\begin{tabular}{cccccc}
\toprule
Corrida $k$ & Empleados & Tasa de empleo & En organización & $\sum$ meses & Tiempo medio \\
\midrule
1 & 86 & $86/100 = 0.86$ & 72 & 415.5 & 5.771 \\
2 & 82 & $82/100 = 0.82$ & 59 & 307.5 & 5.212 \\
3 & 86 & $86/100 = 0.86$ & 61 & 335.0 & 5.492 \\
4 & 89 & $89/100 = 0.89$ & 79 & 458.0 & 5.797 \\
5 & 83 & $83/100 = 0.83$ & 71 & 327.5 & 4.613 \\
6 & 87 & $87/100 = 0.87$ & 70 & 295.5 & 4.221 \\
7 & 89 & $89/100 = 0.89$ & 71 & 339.5 & 4.782 \\
8 & 90 & $90/100 = 0.90$ & 77 & 481.0 & 6.247 \\
9 & 88 & $88/100 = 0.88$ & 69 & 386.5 & 5.601 \\
10 & 88 & $88/100 = 0.88$ & 66 & 344.5 & 5.220 \\
11 & 89 & $89/100 = 0.89$ & 69 & 427.0 & 6.188 \\
12 & 89 & $89/100 = 0.89$ & 69 & 468.0 & 6.783 \\
13 & 90 & $90/100 = 0.90$ & 78 & 401.0 & 5.141 \\
14 & 87 & $87/100 = 0.87$ & 75 & 389.0 & 5.187 \\
15 & 87 & $87/100 = 0.87$ & 61 & 332.5 & 5.451 \\
16 & 84 & $84/100 = 0.84$ & 61 & 284.0 & 4.656 \\
17 & 87 & $87/100 = 0.87$ & 70 & 384.5 & 5.493 \\
18 & 86 & $86/100 = 0.86$ & 71 & 387.0 & 5.451 \\
19 & 89 & $89/100 = 0.89$ & 70 & 390.0 & 5.571 \\
20 & 87 & $87/100 = 0.87$ & 70 & 393.5 & 5.621 \\
21 & 88 & $88/100 = 0.88$ & 68 & 354.0 & 5.206 \\
22 & 87 & $87/100 = 0.87$ & 66 & 337.0 & 5.106 \\
23 & 87 & $87/100 = 0.87$ & 72 & 390.0 & 5.417 \\
24 & 88 & $88/100 = 0.88$ & 72 & 419.0 & 5.819 \\
25 & 88 & $88/100 = 0.88$ & 70 & 373.5 & 5.336 \\
\bottomrule
\end{tabular}
\end{mitabla}

Con las $n = 1000$ corridas se calcula la media, la desviación estándar y el intervalo de confianza de cada indicador, con el mismo procedimiento de los Pasos 2 a 6 del intervalo de confianza ($n - 1 > 30$, por lo que $t \approx 1.96$):

\begin{mitabla}{Resumen de las 1000 corridas por escenario \datosprueba}
\begin{tabular}{llcccc}
\toprule
Indicador & Escenario & $\bar{x}$ & $S$ & $\frac{S}{\sqrt{n}} \cdot 1.96$ & IC 95\% \\
\midrule
Tasa de empleo & Base & 0.8722 & 0.0334 & 0.0021 & $[0.8701,\ 0.8742]$ \\
 & Modificado & 0.9353 & 0.0248 & 0.0015 & $[0.9338,\ 0.9369]$ \\
Tiempo medio (meses) & Base & 5.348 & 0.541 & 0.034 & $[5.315,\ 5.382]$ \\
 & Modificado & 4.315 & 0.478 & 0.030 & $[4.285,\ 4.344]$ \\
\bottomrule
\end{tabular}
\end{mitabla}

\textit{Validación del modelo.} La media simulada del escenario base (tasa 87.2\%, tiempo 5.35 meses) coincide con lo observado en la encuesta (87.0\% y 5.32 meses), lo que indica que el modelo reproduce los datos de los que parte.

\textit{Comparación de escenarios.} Diferencia en la tasa de empleo: $0.9353 - 0.8722 = 0.0632$ (6.3 puntos porcentuales). Diferencia en el tiempo medio: $4.315 - 5.348 = -1.034$ meses. \datosprueba



\subsection{Respuesta a los Criterios ARCU-SUR}\label{sec:criterios}

La Guía de Autoevaluación ARCU-SUR pide, para cada componente, analizar
la carrera en relación con sus criterios y mencionar los aspectos
favorables y desfavorables (ARCU-SUR, 2016a). A continuación se resume,
para cada criterio al que este estudio aporta evidencia, qué pide el
criterio, qué resultado de este trabajo lo respalda y cómo se valora.
Se considera favorable un indicador cuando la mayoría de los
encuestados (más del 50\%) responde en el sentido que el criterio
espera, y desfavorable en caso contrario; los indicadores que no tienen
un valor de referencia se reportan como línea base.

% Generado por sistema/backend ReporteTex.java -- no editar a mano.
% Datos: Base de datos: 148 importado (titulados) / Base de datos: 53 importado (empleadores)
\paragraph{3.2.3 Condiciones de empleo.} \textbf{Qué pide:} que el diseño de la carrera y la calidad de la formación se reflejen en las condiciones de empleo de los graduados; sus indicadores son el tiempo medio para obtener el primer empleo, el destino laboral y la tasa de empleo, y la concordancia entre titulación y empleo. \textbf{Evidencia del estudio:} tasa de empleo de 83.1\% (123 de 148, Sección~\ref{sec:titulados}); tiempo medio para obtener el primer empleo de 2.63 meses, con IC 95\% $[1.85,\ 3.42]$; el 80.6\% de quienes trabajan en una organización lo hace en el sector privado y el rubro más frecuente es ``Desarrollo de Software''; cruces C2 (sin relación significativa, $X_0^2 = 0.002$ frente a $3.841$) y C6 (sin relación significativa, $X_0^2 = 0.227$ frente a $3.841$). El tiempo medio no tiene un valor de referencia en ARCU-SUR, por lo que queda como línea base para el seguimiento.

\textbf{Valoración.} Aspectos favorables: titulados empleados (83.1\%); titulados que consideran pertinentes a su formación las labores de su puesto (90.7\%); titulados que perciben concordancia entre su formación y los requerimientos del mercado (61.9\%). No se identifican aspectos desfavorables.

\paragraph{3.2.1 Resultados.} \textbf{Qué pide:} evaluar el resultado del proceso formativo y utilizar esa evaluación para realizar ajustes. \textbf{Evidencia del estudio:} el 74.1\% de los titulados está de acuerdo con estar satisfecho con la formación recibida (media 2.78 en la escala de 1 a 4); el aspecto a mejorar más señalado es ``Actualizar el plan de estudios'' (71.6\%); cruce C1 (sin relación significativa, $X_0^2 = 0.796$ frente a $5.991$).

\textbf{Valoración.} Aspectos favorables: titulados satisfechos con la formación (74.1\%). No se identifican aspectos desfavorables.

\paragraph{3.2.2 Vinculación y seguimiento a los graduados.} \textbf{Qué pide:} un sistema de seguimiento de los graduados que permita conocer sus condiciones de empleo, e instancias de participación de los graduados para mejorar la carrera. \textbf{Evidencia del estudio:} la encuesta aplicada y la herramienta desarrollada constituyen un mecanismo de seguimiento: cada nueva aplicación se importa y los indicadores se recalculan. El 57.4\% de los titulados manifiesta interés en formar parte de una red oficial de contactos de la Carrera, lo que da una base para esa participación.

\textbf{Valoración.} Aspectos favorables: titulados interesados en la red de contactos (57.4\%). No se identifican aspectos desfavorables.

\paragraph{2.1.2 Perfil de egreso.} \textbf{Qué pide:} coherencia entre el perfil de egreso y la demanda de competencias profesionales expresada por agentes sociales relevantes. \textbf{Evidencia del estudio:} valoración de los empleadores en la escala de perfil profesional (Sección~\ref{sec:empleadores}), donde se considera de acuerdo la suma de ``totalmente'' y ``parcialmente de acuerdo''.

\textbf{Valoración.} Aspectos favorables: ``El título es consistente con las competencias'' (86.8\%); ``Conozco el perfil profesional declarado'' (84.9\%); ``El perfil es coherente con el desempeño'' (79.2\%); ``Evidencian competencias laborales'' (86.8\%); ``Demuestran valores y ética profesional'' (83.0\%). No se identifican aspectos desfavorables.

\paragraph{2.1.7 Actualización curricular.} \textbf{Qué pide:} mecanismos permanentes de actualización del currículo, en consonancia con los desarrollos profesionales. \textbf{Evidencia del estudio:} la asignatura más señalada como ventaja competitiva es ``Base de Datos'' (94.6\%) y la más señalada como poco útil es ``Aplic. Interactivas TV Digital'' (48.6\%); el 71.6\% señala la actualización del plan de estudios como aspecto a mejorar. Las tecnologías y habilidades mencionadas por los empleadores (Sección~\ref{sec:empleadores}) son un insumo directo para esa actualización.

\paragraph{2.4.1 Cursos de actualización profesional y 1.1.5 Posgrado.} \textbf{Qué piden:} extender los conocimientos hacia los graduados mediante cursos de actualización y programas de posgrado afines a la carrera, y promover estudios de posgrado. \textbf{Evidencia del estudio:} el 66.2\% tiene o cursa formación complementaria y el 83.1\% está interesado en un posgrado del área; el área más solicitada es ``IA''; de los interesados, el 41.5\% elegiría la UMSS (EUPG o FCyT); cruces C3 (sin relación significativa, $X_0^2 = 0.769$ frente a $3.841$), C4 (sin relación significativa, $X_0^2 = 1.749$ frente a $3.841$) y C5 (sin relación significativa, $X_0^2 = 2.013$ frente a $9.488$).

\textbf{Valoración.} Aspectos favorables: titulados interesados en un posgrado (83.1\%); titulados que consideran afín la oferta de posgrado de la FCyT (70.9\%). Aspectos desfavorables: interesados que elegirían la UMSS para su posgrado (41.5\%).

\paragraph{2.4.2 Relaciones con el sector público y privado.} \textbf{Qué pide:} que la institución establezca relaciones con empresas y organizaciones públicas y privadas para cooperar en actividades conjuntas. \textbf{Evidencia del estudio:} encuesta a organizaciones de ambos sectores (Sección~\ref{sec:empleadores}).

\textbf{Valoración.} Aspectos favorables: organizaciones que contrataron titulados en los últimos 5 años (79.2\%); ``Consideran nuestra opinión en prácticas y pasantías'' (58.5\%); ``He participado en procesos de retroalimentación'' (52.8\%); ``La Carrera mantiene vínculos que generan oportunidades'' (62.3\%); ``He participado en vinculación entre profesionales'' (50.9\%); ``Recurrimos a la UMSS cuando requerimos profesionales'' (60.4\%). Aspectos desfavorables: ``Las autoridades consultan nuestra opinión'' (47.2\%); ``He participado en vinculación con el entorno laboral'' (47.2\%); ``He participado en difusión de oportunidades laborales'' (47.2\%); ``He sido consultado sobre el desempeño de los titulados'' (47.2\%); ``Han consultado nuestra opinión sobre la pertinencia'' (43.4\%); ``He sido consultado en la autoevaluación de la Carrera'' (35.8\%).

\paragraph{1.2.2 Sistemas de información para la toma de decisiones.} \textbf{Qué pide:} sistemas con información relevante, confiable y actualizada que respalden las decisiones institucionales. \textbf{Evidencia del estudio:} la herramienta desarrollada importa las respuestas, elimina los datos personales, calcula los indicadores y cruces de este informe y exporta los resultados, por lo que puede usarse de forma permanente para actualizar esta información.



\section{Conclusiones y Propuestas Estratégicas}

\textit{[Pendiente: las conclusiones se redactan con los resultados
reales. Las líneas de acción de abajo son las que pide el enunciado,
asociadas al indicador que las sustentaría.]}

\begin{itemize}
  \item \textbf{Fomento de becas y alianzas de financiamiento}, si la
  proporción que financia su formación con recursos propios es alta
  (HU3, cruce C7).
  \item \textbf{Posgrados modulares en las áreas de mayor interés},
  priorizando aquellas en que la mayoría elige instituciones distintas
  de la UMSS (U8, cruce C5, criterio 2.4.1).
  \item \textbf{Vinculación empresarial}: pasantías y consulta periódica
  a empleadores, si los criterios ``las autoridades consultan nuestra
  opinión'' y ``consideran nuestra opinión en prácticas'' tienen bajo
  acuerdo (U11, criterios 2.1.2 y 2.4.2), y si la concordancia
  formación--mercado se asocia con un ingreso más rápido al empleo (C6).
  \item \textbf{Observatorio permanente de titulados}: aplicar la encuesta
  de forma periódica y procesarla con la herramienta desarrollada,
  aprovechando el interés en la red de contactos (U10, criterio 3.2.2).
\end{itemize}

\subsection*{Limitaciones}

\begin{itemize}
  \item El formulario no pregunta el año de titulación, por lo que no se
  pueden analizar resultados por cohorte.
  \item El tiempo al primer empleo solo lo responde quien hoy trabaja en
  una organización, y su media es una cota inferior.
  \item La escala de empleadores no es simétrica y algunos de sus ítems
  agrupan varias afirmaciones.
  \item Con cinco organizaciones, el análisis de empleadores es solo
  descriptivo.
\end{itemize}

\section*{Declaración de Uso de Inteligencia Artificial}
\addcontentsline{toc}{section}{Declaración de Uso de Inteligencia Artificial}

En la elaboración de este trabajo se utilizó el asistente de
inteligencia artificial Claude (Anthropic, 2026), mediante la
herramienta Claude Code, entre septiembre y octubre de 2026. La
Tabla~\ref{tab:ia} detalla en qué tareas se usó y qué parte quedó a
cargo del autor.

\begin{mitabla}{Tareas en las que se usó inteligencia artificial\label{tab:ia}}
\begin{tabular}{L{3.6cm}L{5.6cm}L{5.4cm}}
\toprule
Tarea & Uso de la IA & Revisión y decisiones del autor \\
\midrule
Material de referencia & Transcripción a LaTeX de los capítulos de Coss Bu y de los apuntes de clase a partir de los PDF, marcando los errores del original con notas al pie & Revisión de las transcripciones \textit{[Completar]} \\
Criterios ARCU-SUR & Lectura de los documentos de Criterios, Autoevaluación y Pares, y propuesta de qué preguntas del formulario responden a cada indicador & Selección final de variables y cruces \textit{[Completar]} \\
Revisión de las encuestas & Identificación de problemas de diseño del formulario y propuesta de decisiones de procesamiento (punto medio de rangos, agrupaciones, escala de empleadores) & Decisión de no modificar el formulario y de enviarlo solo a titulados de la UMSS; aprobación de las decisiones de procesamiento \\
Herramienta de software & Generación del código del backend (Java) y del frontend (Vue, Naive UI, ECharts) & Pruebas de funcionamiento y ajustes \textit{[Completar]} \\
Dataset de prueba & Generación de respuestas sintéticas para validar la herramienta & --- (no son resultados del estudio; están marcados en el informe) \\
Modelo de simulación & Propuesta del modelo (generador congruencial mixto y Transformada Inversa discreta) y del escenario modificado & Validación del modelo según los métodos de la asignatura \textit{[Completar]} \\
Redacción del informe & Borrador de las secciones 1 a 6 y de las tablas & Revisión y corrección del texto \textit{[Completar]} \\
\bottomrule
\end{tabular}
\end{mitabla}

No se usó inteligencia artificial para generar respuestas de las
encuestas presentadas como reales: los resultados finales provienen
exclusivamente de los datos recolectados y procesados con la
herramienta. Tampoco se usó para \textit{[Completar: por ejemplo,
aplicación de las encuestas, contacto con las organizaciones, reuniones
con el docente]}.

\section*{Referencias Bibliográficas}
\addcontentsline{toc}{section}{Referencias Bibliográficas}

\begin{hangparas}{1.27cm}{1}
Ali Rodriguez, J. G. (2021). \textit{Simulación de Sistemas: apuntes de
clases} [Apuntes de la materia dictada por H. F. Villarroel Tapia].
Carrera de Ingeniería de Sistemas, Universidad Mayor de San Simón.

Anthropic. (2026). \textit{Claude} (versión Opus 5.5) [Modelo de
lenguaje de gran tamaño]. \url{https://claude.ai}

ARCU-SUR. (2016a). \textit{Guía de autoevaluación para carreras de
Ingeniería}. Red de Agencias Nacionales de Acreditación, Sector Educativo
del MERCOSUR.
\url{https://www.arcusur.org/wp-content/uploads/2024/02/Guia_de_Autoevaluacion_Ingenieria.pdf}

ARCU-SUR. (2016b). \textit{Guía de evaluación para los comités de pares
evaluadores: Titulación Ingeniería}. Red de Agencias Nacionales de
Acreditación, Sector Educativo del MERCOSUR.
\url{https://www.arcusur.org/wp-content/uploads/2024/02/Guia_de_Pares_Ingenieria.pdf}

ARCU-SUR. (2019). \textit{Criterios de calidad para la acreditación
ARCU-SUR: Ingeniería}. Red de Agencias Nacionales de Acreditación,
Sector Educativo del MERCOSUR.
\url{https://www.arcusur.org/wp-content/uploads/2024/03/Criterios_de_Ingenieria_Mayo_2019.pdf}

Coss Bu, R. (1998). \textit{Simulación: un enfoque práctico}. Limusa.
\end{hangparas}

\appendix

\section{Tabla de Homologación de Cargos}\label{anexo:homologacion}

\begin{mitabla}{Homologación de los cargos de ambas encuestas}
\begin{tabular}{lL{5.2cm}L{5.8cm}}
\toprule
Grupo & Encuesta a titulados & Encuesta a empleadores \\
\midrule
Soporte                 & Help Desk            & De apoyo / asistencia tecnológica \\
Técnico / Desarrollo    & Técnico, Desarrollo  & Técnico \\
Analista / Especialista & ---                  & Analista, especialista, consultor tecnológico \\
Gestión / Jefatura      & Ejecutivo            & Supervisor / coordinador, jefatura o gerencial \\
Otro                    & Otro                 & Otro \\
\bottomrule
\end{tabular}
\end{mitabla}

\end{document}
